\chapter{PRELIMINARIES AND MANIFOLDS}

\textit{July 20th.}

The primary reference for this course will be Loring Tu's \emph{An Introduction to Manifolds}. Prerequisites include Analysis of Several Variables, Linear Algebra, and Topology. The mid-term will account for 35\% of the final grade, quizzes and student lectures will account for 15\%, and the final exam will account for 50\%. 

Simply put, a manifold is any object that is locally Euclidean.

\section{Calculus on $\R^{n}$}
A point $p$ or element in $\R^{n}$ is associated with an $n$-tuple of coordinates $x^{1},x^{2},\ldots,x^{n}$. Let $U$ be an open neighbourhood of $p$ in $\R^{n}$. For a non-negative integer $k$, we say that a function $f:U\to\R$ is \eax{$C^{k}$-continuous} at $p \in U$ if all partial derivatives of $f$ of order $\leq k$ exist and are continuous at $p$. $f$ is said to be $C^{\infty}$ or \eax{smooth} if $f$ is $C^{k}$ for all $k \geq 0$. $f$ is $C^{k}$ (or $C^{\infty}$)on $U$ if it is $C^{k}$ (or $C^{\infty}$) at every point in $U$.

If $f$ happens to be a vector-valued function, i.e. $f:U\to\R^{m}$, we say $f = (f^{1},f^{2},\ldots,f^{m})$ is $C^{k}$ (or $C^{\infty}$) at $p$ if each component function $f^{i}:U\to\R$ is $C^{k}$ (or $C^{\infty}$) at $p$. $f$ is $C^{k}$ (or $C^{\infty}$) on $U$ if it is $C^{k}$ (or $C^{\infty}$) at every point in $U$.

\begin{definition}
    Let $R$ be a commutative ring with unity. We term $A$ to be an $R$-\eax{algebra} if $A$ is a commutative ring with unity and an $R$-module, and for $r \in R$ and $a,b \in A$, we have $r(ab) = (ra)b = a(rb)$.
\end{definition}
In this course, we will be primarily be concerned with the case when $R$ is a field, and $A$ is an $R$-vector space.

For a subset $U \subseteq \R^{n}$, we denote $C^{k}(U)$ (or $C^{\infty}(U)$) as the set of all $C^{k}$ (or $C^{\infty}$) functions from $U$ to $\R$. Note that this is a commutative ring with unity. It is also a vector space over $\R$. Together, $C^{k}(U)$ (or $C^{\infty}(U)$) is a (commutative) algebra over $\R$. Moreover, we have the descending chain of subalgebras
\begin{align}
    C^{0}(U) \supseteq C^{1}(U) \supseteq C^{2}(U) \supseteq \cdots C^{k}(U) \supseteq C^{k+1}(U) \supseteq \cdots
\end{align}
with
\begin{align}
    \bigcap_{k \geq 0} C^{k}(U) = C^{\infty}(U).
\end{align}

$f:U \to \R$ is termed \eax{real analytic} at $p \in U$ if there exists a neighbourhood $V \subseteq U$ of $p$ such that
\begin{align}
    f(x) = f(p) + \sum_{i=1}^{n} \frac{\partial f}{\partial x^{i}}(p)(x^{i}-p^{i}) + \cdots + \frac{1}{k!} \sum_{1 \leq i_{1},\ldots,i_{k} \leq n} \frac{\partial^{k} f}{\partial x^{i_{1}} \cdots \partial x^{i_{k}}}(p)(x^{i_{1}}-p^{i_{1}}) \cdots (x^{i_{k}}-p^{i_{k}}) + \cdots
\end{align}
for all $x \in V$. $f$ is said to be real analytic on $U$ if it is real analytic at every point in $U$. We denote the set of all real analytic functions on $U$ as $C^{\omega}(U)$. Note that $C^{\omega}(U) \subseteq C^{\infty}(U)$.

\begin{theorem}[\eax{Taylor's theorem}]
    Let $p \in U \subseteq \R^{n}$, and $f:U \to \R$ be a $C^{\infty}$ function. Further assume that for all $x \in U$, the line segment from $p$ to $x$ is contained in $U$ ($U$ in this case is said to be \eax{star-shaped} with respect to $p$). Then there exist $g_{1},\ldots,g_{n} \in C^{\infty}(U)$ such that, for all $x \in U$,
    \begin{align}
        f(x) = f(p) + \sum_{i=1}^{n} g_{i}(x)(x^{i}-p^{i}),
    \end{align}
    with
    \begin{align}
        g_{i}(p) = \frac{\partial f}{\partial x^{i}}(p), \quad 1 \leq i \leq n.
    \end{align}
\end{theorem}
\begin{proof}
    Let $x \in U$. Then the line segment $\{p+t(x-p) \mid t \in [0,1]\}$ is contained in $U$. Composing $f$ with the map $t \mapsto p+t(x-p)$, we obtain
    \begin{align}
        \frac{df}{dt}(p+t(x-p)) = \sum_{i=1}^{n} \frac{\partial f}{\partial x^{i}}(p+t(x-p))(x^{i}-p^{i}).
    \end{align}
    Integrating from $0$ to $1$, we get
    \begin{align}
        f(x) - f(p) = \int_{0}^{1} \frac{df}{dt}(p+t(x-p)) dt = \sum_{i=1}^{n} (x^{i}-p^{i}) \int_{0}^{1} \frac{\partial f}{\partial x^{i}}(p+t(x-p)) dt = \sum_{i=1}^{n} (x^{i}-p^{i}) g_{i}(x).
    \end{align}
    Clearly, because $f$ is $C^{\infty}$, each $g_{i}$ is also $C^{\infty}$. Further, we have $g_{i}(p) = \int_{0}^{1} \frac{\partial f}{\partial x^{i}}(p) dt = \frac{\partial f}{\partial x^{i}}(p)$.
\end{proof}

Let $p = (p^{1},\ldots,p^{n}), v \in \R^{n}$, and $f:U \to \R$ be $C^{\infty}$ at $p$. The \eax{directional derivative} of$ f$ at $p$ in the direction of $v$ is defined as
\begin{align}
    \lim_{t \to 0} \frac{f(p+tv)-f(p)}{t} = \frac{d}{dt}f(p+tv)\Big|_{t=0} = \sum_{i=1}^{n} \frac{df}{dx^{i}}(p)v^{i}.
\end{align}
From here, we obtain the operator
\begin{align}
    D_{v} = \sum_{i=1}^{n} v^{i} \frac{\partial}{\partial x^{i}}\Big|_{p}
\end{align}
Now consider all pairs of the form $(f,U)$ where $p \in U \subseteq \R^{n}$ is an open neighbourhood of $p$ and $f:U \to \R$ is $C^{\infty}$. Define an equivalence relation $\sim$ on this set by declaring $(f,U) \sim (g,V)$ if and only if there exists an open neighbourhood $W \subseteq U \cap V$ of $p$ such that $f|_{W} = g|_{W}$. The equivalence class of $(f,U)$ is denoted by $[(f,U)]$, and the set of all such equivalence classes is denoted by $C^{\infty}_{p}(\R^{n})$. Such an equivalence class $[(f,U)]$ is called the \eax{germ} of $f$ at $p$. We can give $C^{\infty}_{p}(\R^{n})$ the structure of a commutative ring with unity by defining
\begin{align}
    [(f,U)] + [(g,V)] = [(f+g,U \cap V)], \quad [(f,U)] \cdot [(g,V)] = [(fg,U \cap V)], \quad 1 = [(1,\R^{n})].
\end{align}
It can also be given the structure of a vector space over $\R$ by defining
\begin{align}
    \lambda [(f,U)] = [(\lambda f,U)], \quad \lambda \in \R.
\end{align}
Thus, $C^{\infty}_{p}(\R^{n})$ is an $\R$-algebra. Now, given a germ $\alpha = [(f,U)] \in C^{\infty}_{p}(\R^{n})$ and a vector $v \in \R^{n}$, we define
\begin{align}
    D_{v}\alpha = \sum_{i=1}^{n} v^{i} \frac{\partial f}{\partial x^{i}}\Big|_{p} = v \cdot \nabla f(p).
\end{align}
Note that the choice of $f$ is immaterial, since if $[(f,U)] = [(g,V)]$, then there exists an open neighbourhood $W \subseteq U \cap V$ of $p$ such that $f|_{W} = g|_{W}$, and hence $\nabla f(p) = \nabla g(p)$.
One can infer that $D_{\lambda v}(f) = \lambda D_{v}(f) = D_{v}(\lambda f)$ and $D_{v}(f+g) = D_{v}(f) + D_{v}(g)$ for all $\lambda \in \R$ and $f,g \in C^{\infty}_{p}(\R^{n})$. Thus, $D_{v}$ is an $\R$-linear map from $C^{\infty}_{p}(\R^{n})$ to $\R$ (hereforth, we shall use $f$ to also denote the germ $[(f,U)]$ of $f$ at $p$). Also,
\begin{align}
    D_{v}(fg) = \sum_{i=1}^{n} v^{i} \frac{\partial (fg)}{\partial x^{i}}(p) = \sum_{i=1}^{n} v^{i} \left( f(p) \frac{\partial g}{\partial x^{i}}(p) + g(p) \frac{\partial f}{\partial x^{i}}(p) \right) = f(p) D_{v}(g) + g(p) D_{v}(f).
\end{align} 
Such a property makes $D_{v}$ what we call a derivation.
\begin{definition}
    Any $\R$-linear map $T:C^{\infty}_{p}(\R^{n}) \to \R$ satisfying the Leibniz rule $T(fg) = f(p)T(g) + g(p)T(f)$ for all $f,g \in C^{\infty}_{p}(\R^{n})$ is called a \eax{derivation} at $p$.
\end{definition}
One may show that if $T$ is a derivation at $p$, then $T(1) = 0$. 

The set of all derivations at $p$ is denoted by $\mf{D}_{p}(\R^{n})$. It is easy to see that $\mf{D}_{p}(\R^{n})$ is a vector space over $\R$. Now consider the map
\begin{align}
    \theta:\R^{n} \to \mf{D}_{p}(\R^{n}), \quad v \mapsto D_{v}.
\end{align}
This $\theta$ is an $\R$-linear map (verify). We claim that $\theta$ is an isomorphism of real vector spaces. To this end, let $v \in \R^{n}$ such that $D_{v} = 0 \in \mf{D}_{p}(\R^{n})$. Then $D_{v}f = 0$ for all $f \in C^{\infty}_{p}(\R^{n})$. In particular, for the coordinate functions $x^{i}$, we have $D_{v}x^{i} = v_{i} = 0$ for all $1 \leq i \leq n$. Thus, $v = 0$, and hence $\theta$ is injective. To show surjectivity, we appeal to Taylor's theorem; let $D \in \mf{D}_{p}(\R^{n})$. We claim that $D = D_{v}$ where $v = (D(x^{1}),\ldots,D(x^{n}))$. Indeed, for any $f \in C^{\infty}_{p}(\R^{n})$, we have
\begin{align}
    D(f) &= D(f(p)) + D\left( \sum_{i=1}^{n} (x^{i}-p^{i}) g_{i}(x) \right) = 0 + \sum_{i=1}^{n} \left( (x^{i}-p^{i})(p) Dg_{i}(x) + D(x^{i}-p^{i})g_{i}(p) \right) \notag \\
    &= \sum_{i=1}^{n} (Dx^{i})g_{i}(p) = \sum_{i=1}^{n} \frac{\partial f}{\partial x^{i}}(p) Dx^{i} = D_{v}(f).
\end{align}
Hence, $\theta$ is surjective, and we have the desired isomorphism $\mf{D}_{p}(\R^{n}) \cong \R^{n}$.

Note that a germ $f \in C^{\infty}_{p}(\R^{n})$ has a valid value at $p$, but no other point. Hence, we have the evaluation map 
\begin{align}
    \ev_p: C^{\infty}_{p}(\R^{n}) \to \R, \quad [(f,U)] \mapsto f(p).
\end{align}
The kernel here is $\{[(f,U)] \in C^{\infty}_{p}(\R^{n}) \mid f(p) = 0\}$, which is a ideal in $C^{\infty}_{p}(\R^{n})$. For any other element $[(f,U)] \in C^{\infty}_{p}(\R^{n})$ with $f(p) \neq 0$, a valid inverse is given by $[(f^{-1},U\setminus\{f=0\})]$. This makes $C^{\infty}_{p}(\R^{n})$ a local ring.
\\

\section{Multilinear Algebra}

\textit{July 22nd.}

Hereforth, $V$ shall denote a finite-dimensional real vector space, and $V^{V}$ shall denote its \eax{dual space} $\Hom_{\R}(V,\R)$. If $e_{1},\ldots,e_{n}$ are a basis for $V$, then the \eax{dual basis} of $V^{V}$ is denoted by $\alpha^{1},\ldots,\alpha^{n} \in V^{V}$, and is defined by $\alpha^{i}(e_{j}) = \delta_{i,j}$ for all $1 \leq i,j \leq n$. One may verify that the dual basis $\{\alpha^{1},\ldots,\alpha^{n}\}$, induced by the basis $\{e_{1},\ldots,e_{n}\}$, is indeed a basis for $V^{V}$. $S_{k}$ shall denote the permutation group on $k$ symbols, satisfying $\abs{S_{k}} = k!$. For any $\sigma \in S_{k}$, the sign of $\sigma$ is denoted by $\sgn(\sigma)$, and is defined as $+1$ if $\sigma$ is an even permutation, and $-1$ if $\sigma$ is an odd permutation. 

With notation out of the way, let $k \geq 1$ be an integer. We shall consider the $k$-tuple $V^{k} = V \times \cdots \times V$ ($k$ times).

\begin{definition}
    A mapping $f:V^{k} \to \R$ is said to be \eax{$k$-linear} if $f$ is linear in each of its $k$ arguments.
\end{definition}
A special case is when $k = 2$, in which case $f$ is called a \eax{bilinear form}. Such an $f:V \times V \to \R$ satisfies
\begin{align}
    f(\lambda v + \mu w, u) = \lambda f(v,u) + \mu f(w,u), \quad f(u,\lambda v + \mu w) = \lambda f(u,v) + \mu f(u,w).
\end{align}
A common example of a bilinear form is any valid inner product $\ip{\cdot,\cdot}:V \times V \to \R$ on $V$. Another example is the determinant $\det:\R^{n} \to \R$, where the vectors are interpreted as the columns of a matrix. We can also write
\begin{align}
    L_{k}(V) = \{f:V^{k} \to \R \mid f \text{ is $k$-linear}\}.
\end{align}
It is easy to see that $L_{k}(V)$ is an $\R$-vector space.

\begin{definition}
    A $k$-linear map $f:V^{k} \to \R$ is called a \eax{symmetric $k$-linear form} if for all $\sigma \in S_{k}$,
    \begin{align}
        f(v_{1},\ldots,v_{k}) = f(v_{\sigma(1)},\ldots,v_{\sigma(k)}).
    \end{align}
    That is, $f$ is invariant under any permutation of its arguments. Moreover, $f$ is called an \eax{alternating $k$-linear form} if for all $\sigma \in S_{k}$,
    \begin{align}
        f(v_{1},\ldots,v_{k}) = \sgn(\sigma)f(v_{\sigma(1)},\ldots,v_{\sigma(k)}).
    \end{align}
\end{definition}
One can see that the standard inner product on $\R^{n}$ is a symmetric bilinear form, and the determinant is an alternating $n$-linear form. Related to these definitions, we can immediately define the following subspaces of $L_{k}(V)$:
\begin{align}
    S_{k}(V) = \{f \in L_{k}(V) \mid f \text{ is symmetric}\} \subseteq L_{k}(V), \quad A_{k}(V) = \{f \in L_{k}(V) \mid f \text{ is alternating}\} \subseteq L_{k}(V).
\end{align}
Currently, our main interest of study is $A_{k}(V)$, the space of alternating $k$-linear forms on $V$. Elements of $A_{k}(V)$ are termed \eax{alternating $k$-tensors}. Similarly, those of $S_{k}(V)$ and $L_{k}(V)$ are called \eax{symmetric $k$-tensors} and \eax{$k$-tensors}, respectively. For $k = 0$, we define $L_{0}(V) = S_{0}(V) = A_{0}(V) = \R$. For $k = 1$, we have $L_{1}(V) = S_{1}(V) = A_{1}(V) = V^{V}$. We can define the (left) action of $S_{k}$ on $L_{k}(V)$ by
\begin{align}
    (\sigma \cdot f)(v_{1},\ldots,v_{k}) = f(v_{\sigma(1)},\ldots,v_{\sigma(k)}).
\end{align}
Thus, $f$ is symmetric if $\sigma \cdot f = f$ for all $\sigma \in S_{k}$, and $f$ is alternating if $\sigma \cdot f = \sgn(\sigma)f$ for all $\sigma \in S_{k}$.

Given an $f \in L_{k}(V)$, we can define the \eax{symmetrization} $Sf \in S_{k}(V)$, of $f$, as
\begin{align}
    Sf = \sum_{\sigma \in S_{k}} \sigma \cdot f.
\end{align}
One may, indeed, check that $Sf$ is symmetric. Similarly, we can define the \eax{alternation} $Af \in A_{k}(V)$, of $f$, as
\begin{align}
    Af = \sum_{\sigma \in S_{k}} \sgn(\sigma) \sigma \cdot f.
\end{align}

\subsection{The Tensor and Wedge Products}

\noindent \textit{July 24th.}

For any $f \in L_{k}(V)$ and $g \in L_{l}(V)$, we can define the \eax{tensor product} $f \otimes g \in L_{k+l}(V)$ as
\begin{align}
    (f \otimes g)(v_{1},\ldots,v_{k+l}) = f(v_{1},\ldots,v_{k})g(v_{k+1},\ldots,v_{k+l}).
\end{align}
Verify that $f \otimes g$ is indeed $(k+l)$-linear. We are interested in the tensor product of alternating tensors, so let $f \in A_{k}(V)$ and $g \in A_{l}(V)$. Then, we can define the \eax{wedge product} $f \wedge g \in A_{k+l}(V)$ as
\begin{align}
    f \wedge g = \frac{1}{k!l!} A(f \otimes g) = \frac{1}{k!l!} \sum_{\sigma \in S_{k+l}} \sgn(\sigma) \sigma \cdot (f \otimes g).
\end{align}
The associativity of the tensor product holds easily: for $f \in L_{k}(V)$, $g \in L_{l}(V)$, and $h \in L_{m}(V)$,
\begin{align}
    (f \otimes g) \otimes h = f \otimes (g \otimes h).
\end{align}

\begin{proposition}
    For $f \in A_{k}(V)$ and $g \in A_{l}(V)$, we have
    \begin{align}
        f \wedge g = (-1)^{kl} g \wedge f.
    \end{align}
\end{proposition}
\begin{proof}
    We shall show that $A(f \otimes g) = (-1)^{kl} A(g \otimes f)$. To this end, fix the permutation $\tau \in S_{k+l}$ defined by
    \begin{align}
        \tau = \begin{pmatrix}
            1 & 2 & \cdots & l & l+1 & l+2 & \cdots & l+k \\
            k+1 & k+2 & \cdots & k+l & 1 & 2 & \cdots & k
        \end{pmatrix}.
    \end{align}
    We have
    \begin{align}
        A(f \otimes g)(v_{1},\ldots,v_{k+l}) &= \sum_{\sigma \in S_{k+l}} \sgn(\sigma) f(v_{\sigma(1)},\ldots,v_{\sigma(k)}) g(v_{\sigma(k+1)},\ldots,v_{\sigma(k+l)}) \notag \\
        &= \sum_{\sigma \in S_{k+l}} \sgn(\sigma) f(v_{\sigma\tau(l+1)},\ldots,v_{\sigma\tau(l+k)}) g(v_{\sigma\tau(1)},\ldots,v_{\sigma\tau(l)}) \notag \\
        &= \sgn(\tau) \sum_{\sigma \in S_{k+l}} \sgn(\sigma\tau) f(v_{\sigma\tau(l+1)},\ldots,v_{\sigma\tau(l+k)}) g(v_{\sigma\tau(1)},\ldots,v_{\sigma\tau(l)}) \notag \\
        &= \sgn(\tau)A(g \otimes f)(v_{1},\ldots,v_{k+l}).
    \end{align}
    With $\sgn(\tau) = (-1)^{kl}$, the proof is complete.
\end{proof}

\begin{corollary}
    For $f \in A_{k}(V)$, we have $f \wedge f = 0$ if $k$ is odd.
\end{corollary}
\begin{proof}
    Easy since $f \wedge f = (-1)^{k^{2}} f \wedge f = -f \wedge f$ if $k$ is odd, so $f \wedge f = 0$.
\end{proof}

\begin{lemma}
    For $f \in L_{k}(V)$ and $g \in L_{l}(V)$, we have
    \begin{align}
        A(Af \otimes g) = k!A(f \otimes g), \quad A(f \otimes Ag) = l!A(f \otimes g).
    \end{align}
\end{lemma}

\begin{proof}
    We shall prove the first equality; the second follows similarly. We have
    \begin{align}
        A(Af \otimes g) &= \sum_{\sigma \in S_{k+l}} \sgn(\sigma) \sigma \cdot \left( \sum_{\tau \in S_{k}} \sgn(\tau) (\tau \cdot f)  \otimes g \right) = \sum_{\sigma \in S_{k+l}} \sgn(\sigma) \sigma \cdot \left( \sum_{\tau \in S_{k}} \sgn(\tau) \tau \cdot(f \otimes g) \right) \notag \\
        &= \sum_{\tau \in S_{k}} \sum_{\sigma \in S_{k+l}} \sgn(\sigma \tau) \sigma \tau \cdot (f \otimes g) = k! \sum_{\alpha \in S_{k+l}} \sgn(\alpha) \alpha \cdot (f \otimes g) = k! A(f \otimes g).
    \end{align}
\end{proof}

\begin{proposition}
    The wedge product is associative, i.e. for $f \in A_{k}(V)$, $g \in A_{l}(V)$, and $h \in A_{m}(V)$, we have
    \begin{align}
        (f \wedge g) \wedge h = f \wedge (g \wedge h).
    \end{align}
\end{proposition}
\begin{proof}
    We have
    \begin{align}
        (f \wedge g) \wedge h &= \frac{1}{(k+l)!m!} A((f \wedge g) \otimes h) = \frac{1}{(k+l)!m!k!l!} A(A(f \otimes g) \otimes h) = \frac{(k+l)!}{(k+l)!l!k!m!}A((f \otimes g) \otimes h) \notag \\
        &= \frac{1}{k!l!m!} A(f \otimes g \otimes h) = f \wedge (g \wedge h).
    \end{align}
\end{proof}

\begin{corollary}
    $f \wedge g \wedge h = \frac{1}{k!l!m!} A(f \otimes g \otimes h)$ for $f \in A_{k}(V)$, $g \in A_{l}(V)$, and $h \in A_{m}(V)$. Inductively,
    \begin{align}
        f_{1} \wedge \cdots \wedge f_{k} = \frac{1}{k_{1}!\cdots k_{k}!} A(f_{1} \otimes \cdots \otimes f_{k}),
    \end{align}
    where $f_{i} \in A_{k_{i}}(V)$ for $1 \leq i \leq k$.
\end{corollary}

\begin{proposition}
    Let $\alpha^{1},\ldots,\alpha^{n} \in V^{V} = A_{1}(V) = L_{1}(V)$, and $v_{1},\ldots,v_{k} \in V$. Then
    \begin{align}
        (\alpha^{1} \wedge \cdots \wedge \alpha^{k})(v_{1},\ldots,v_{k}) = \det(\alpha^{i}(v_{j}))_{1 \leq i,j \leq k}.
    \end{align}
\end{proposition}

\begin{proof}
    We have
    \begin{align}
        (\alpha^{1} \wedge \cdots \wedge \alpha^{k})(v_{1},\ldots,v_{k}) &= \sum_{\sigma \in S_{k}} \sgn(\sigma)(\alpha^{1} \otimes \cdots \otimes \alpha^{k})(v_{\sigma(1)},\ldots,v_{\sigma(k)}) \notag \\
        &= \sum_{\sigma \in S_{k}} \sgn(\sigma) \alpha^{1}(v_{\sigma(1)})\alpha^{2}(v{\sigma(2)}) \cdots \alpha^{k}(v_{\sigma(k)}) = \det(\alpha^{i}(v_{j}))_{1 \leq i,j \leq k}.
    \end{align}
\end{proof}

\noindent \textit{July 27th.}

Let $V$ be an $\R$-vector space with basis $e_{1},\ldots,e_{n}$, and $V^{V}$ be its dual space with dual basis $\alpha^{1},\ldots,\alpha^{n}$. Moreover, let $I$ be a multi-index of length $k$, i.e. $I = (i_{1},\ldots,i_{k})$ with $1 \leq i_{1},\ldots,i_{k} \leq n$. We define $e_{I} \defeq (e_{i_{1}},\ldots,e_{i_{k}}) \in V^{k}$ and $\alpha^{I} \defeq \alpha^{i_{1}} \wedge \cdots \wedge \alpha^{i_{k}} \in A_{k}(V)$. Then, a $k$-linear form $f \in L_{k}(V)$ is completely determined by its values on the $e_{I}$, and an alternating $k$-linear form $f \in A_{k}(V)$ is completely determined by its values on the $e_{I}$ with $i_{1} < i_{2} < \cdots < i_{k}$.

To this end, we first claim that $\alpha^{I}(e_{J}) = \delta_{J}^{I}$ for all multi-indices $I,J$ of length $k$. Indeed, if $J = (j_{1},\ldots,j_{k})$ with $j_{1} < j_{2} < \cdots < j_{k}$, then
\begin{align}
    \alpha^{I}(e_{J}) = \det(\alpha^{i_{m}}(e_{j_{n}}))_{1 \leq m,n \leq k};
\end{align}.
this is $1$ for $I = J$. Suppose $I \neq J$. Let $1 \leq l \leq k$ be the first index where $i_{l} \neq j_{l}$. Then the $l^{\text{th}}$ row of $\alpha^{i_{m}}(e_{j_{n}})_{1 \leq m,n \leq k}$ is an all zeroes row, and hence $\det(\alpha^{i_{m}}(e_{j_{n}}))_{1 \leq m,n \leq k} = 0$. Thus, $\alpha^{I}(e_{J}) = \delta_{J}^{I}$. We now claim that $\alpha^{I}$, for strictly increasing multi-indices $I$ of length $k$, form a basis for $A_{k}(V)$. For linear independence, we simply have (the sums in this context are over strictly increasing multi-indices $I$ of length $k$):
\begin{align}
    \sum_{I} c_{I} \alpha^{I} = 0 \implies \sum_{I} c_{I} \alpha^{I}(e_{J}) = 0 \implies c_{J} = 0, \quad \forall J.
\end{align}
To show that they span $A_{k}(V)$, let $f \in A_{k}(V)$. Then consider the $k$-linear form
\begin{align}
    g = \sum_{I} f(e_{I}) \alpha^{I}.
\end{align}
We show that $f=g$; it is enough to prove that $f(e_{J}) = g(e_{J})$ for all strictly increasing multi-indices $J$ of length $k$. We have
\begin{align}
    g(e_{J}) = \sum_{I} f(e_{I}) \alpha^{I}(e_{J}) = f(e_{J}) \alpha^{J}(e_{J}) + \sum_{I \neq J} f(e_{I}) \alpha^{I}(e_{J}) = f(e_{J}).
\end{align}

\begin{corollary}
    $A_{k}(V) = 0$ if $k > n$. Further, $\dim A_{k}(V) = \binom{n}{k}$ if $1 \leq k \leq n$.
\end{corollary}


\section{Building Blocks for Manifolds}
Recall that we had defined the class of germs $C_{p}^{\infty}(\R^{n})$ of $C^{\infty}$ functions at a point $p \in \R^{n}$, which agree on some neighbourhood of $p$. One can show that $C_{p}^{\infty}(\R^{n})$ and $C_{p}^{\infty}(U)$ are isomorphic as $\R$-algebras, where $U \subseteq \R^{n}$ is an open neighbourhood of $p$. Note that we term $T_{p}(U) = T_{p}(\R^{n}) = \R^{n}$ as the \eax{tangent space} of $U$ (or $\R$) at $p$. 

\begin{definition}
    Let $U \subseteq \R^{n}$ be an open subset. A \eax{vector field} on $U$ is a function $X:U \to \amalg_{p \in U} T_{p}(\R^{n})$ such that $X(p) \in T_{p}(\R^{n})$ for all $p \in U$. In particular,
    \begin{align}
        X(p) = \sum_{i=1}^{n} a^{i}(p) \frac{\partial}{\partial x^{i}}\Big|_{p}, \quad a^{i}(p) \in \R.
    \end{align}
    We can rewrite this as
    \begin{align}
        X = \sum_{i=1}^{n} a^{i} \frac{\partial}{\partial x^{i}}, \quad a^{i}:U \to \R.
    \end{align}
    $X$ is called a \eax{smooth vector field} if each $a^{i}$ is a smooth function on $U$. The set of all smooth vector fields on $U$ is denoted by $\mf{X}(U)$.
\end{definition}
This is a real vector space; note that for $X,Y \in \mf{X}(U)$, we have $X+Y \in \mf{X}(U)$, and for $\lambda \in \R$, we have $\lambda X \in \mf{X}(U)$. Further, $\mf{X}(U)$ is a $C^{\infty}(U)$-module since for $f \in C^{\infty}(U)$ and $X \in \mf{X}(U)$, we have $fX \in \mf{X}(U)$, where
\begin{align}
    (fX)(p) = f(p)X(p), \quad p \in U.
\end{align}
Notice that $X \in \mf{X}(U)$ describes a map $C^{\infty}(U) \to C^{\infty}(U)$, where $f \mapsto Xf$, being defined as $(Xf)(p)=X_{p}f$. This is a derivation on $C^{\infty}(U)$, since for $f,g \in C^{\infty}(U)$, we have
\begin{align}
    (X(fg))(p) = X_{p}(fg) = f(p)X_{p}g + g(p)X_{p}f = (fXg)(p) + (gXf)(p).
\end{align}
We refine our definition of a derivation as follows.
\begin{definition}
    Let $K$ be a field, and $A$ be a $K$-algebra. A \eax{derivation} $D:A \to A$ is an $\R$-linear map satisfying the Leibniz rule $D(ab) = aD(b) + bD(a)$ for all $a,b \in A$.
\end{definition}

\begin{definition}
    A \eax{differential 1-form} on $U$ is a smooth function $\omega:U \to \amalg_{p \in U} T_{p}(\R^{n})^{V}$ such that $\omega(p) \in T_{p}(\R^{n})^{V} =\vcentcolon T_{p}^{\ast}(\R^{n})$ for all $p \in U$. 
\end{definition}
For example, for $f \in C^{\infty}(U)$, we can define the differential 1-form $df$ on $U$ by $df(p) = df_{p} \in \Hom_{\R}(T_{p}\R^{n},\R)$, where $df_{p}(X) = X_{p}f$ for all $X \in T_{p}\R^{n}$. In particular, $dx^{i}$ is a differential 1-form on $U$ for all $1 \leq i \leq n$. In fact, $\{dx^{1},\ldots,dx^{n}\}$ is the dual basis of $T_{p}^{\ast}\R^{n}$ corresponding to the basis $\{\frac{\partial}{\partial x^{1}}\Big|_{p},\ldots,\frac{\partial}{\partial x^{n}}\Big|_{p}\}$ of $T_{p}\R^{n}$. Thus, a differential 1-form $\omega$ on $U$ can be expressed as
\begin{align}
    \omega(p) = \sum_{i=1}^{n} a_{i}(p) dx^{i}\Big|_{p}, \quad a_{i}:U \to \R.
\end{align}
$T_{p}^{\ast}(\R^{n})$, here, is termed the \eax{cotangent space} of $\R^{n}$ at $p$. Thus, our current setup is as follows: $V = T_{p}\R^{n}$, and $V^{V} = T_{p}^{\ast}\R^{n}$. The basis for $T_{p}\R^{n}$ considered is $\{\frac{\partial}{\partial x^{1}}\Big|_{p},\ldots,\frac{\partial}{\partial x^{n}}\Big|_{p}\}$, and the corresponding dual basis for $T_{p}^{\ast}\R^{n}$ is $\{dx^{1}\Big|_{p},\ldots,dx^{n}\Big|_{p}\}$. And finally, we have
\begin{align}
    A_{k}(T_{p}\R^{n}) = \text{the space of alternating $k$-linear forms on $T_{p}\R^{n}$}
\end{align}
which has the basis $\{dx^{i_{1}}\Big|_{p} \wedge \cdots \wedge dx^{i_{k}}\Big|_{p} \mid 1 \leq i_{1} < \cdots < i_{k} \leq n\}$.

\begin{definition}
    A \eax{differential $k$-form} on $U$ is a smooth function $\omega:U \to \amalg_{p \in U} A_{k}(T_{p}\R^{n})$ such that $\omega(p) \in A_{k}(T_{p}\R^{n})$ for all $p \in U$.
\end{definition}
Any such $\omega$ may be written as $\omega = \sum_{I} a_{I} dx_{I}$. It termed a smooth differential $k$-form if each $a_{I}$ is a smooth function on $U$. The set of all smooth differential $1$-forms on $U$ is denoted by $\Omega^{1}(U)$, and the set of all smooth differential $k$-forms on $U$ is denoted by $\Omega^{k}(U)$. Both these sets are $\R$-vector spaces, and $\Omega^{k}(U)$ is a free $C^{\infty}(U)$-module of rank $\binom{n}{k}$. The isomorphism may be given as
\begin{align}
    \Omega^{1}(U) \cong C^{\infty}(U)^{\oplus \binom{n}{k}}, \quad \omega = \sum_{I} a_{I} dx_{I} \mapsto (a_{1},a_{2},\ldots,a_{\binom{n}{k}}).
\end{align}
We can endow $\Omega^{k}(U)$ with the wedge product $\wedge:\Omega^{k}(U) \times \Omega^{l}(U) \to \Omega^{k+l}(U)$, defined as
\begin{align}
    (\omega \wedge \eta)(p) = \omega(p) \wedge \eta(p), \quad p \in U,
\end{align}
or
\begin{align}
    \omega \wedge \eta = \sum_{I,J} a_{I}b_{J} dx_{I} \wedge dx_{J}
\end{align}
where $I$ and $J$ are strictly increasing multi-indices of length $k$ and $l$, respectively. Note that if $I \cap J \neq \emptyset$, then $dx_{I} \wedge dx_{J} = 0$. Thus, the sum may be taken over all strictly increasing multi-indices $I$ and $J$ of length $k$ and $l$, respectively, such that $I \cap J = \emptyset$. Also,
\begin{align}
    A_{\ast}(V) = \bigoplus_{k=0}^{n} A_{k}(V)
\end{align}
is defined as the \eax{exterior algebra} of $V$. The wedge product on here is associative and anti-commutative. We define $\Omega^{0}(U) = C^{\infty}(U)$. The map $d:\Omega^{k}(U) \to \Omega^{k+1}(U)$ is defined as the \eax{exterior derivative}; for $k = 0$, we have
\begin{align}
    d:C^{\infty}(U) \to \Omega^{1}(U), \quad f \mapsto df = \sum_{i=1}^{n} \frac{\partial f}{\partial x^{i}} dx^{i}.
\end{align}
For a $k$-form $\omega = \sum_{I} a_{I} dx_{I} \in \Omega^{k}(U)$, we have
\begin{align}
    d\omega = \sum_{I} da_{I} \wedge dx_{I} = \sum_{I} \sum_{j=1}^{n} \frac{\partial a_{I}}{\partial x^{j}} dx^{j} \wedge dx_{I}.
\end{align}
\begin{example}
    Let $U \subseteq \R^{3}$. Let us call the coordinates $x^{1} = x$, $x^{2} = y$, and $x^{3} = z$. The only non-zero forms are $\Omega^{0}(U)$, $\Omega^{1}(U)$, $\Omega^{2}(U)$, and $\Omega^{3}(U)$. Note that $\Omega^{0}(U) = C^{\infty}(U)$, and there's nothing to say about this. For $\Omega^{1}(U)$, a form looks like $f\, dx + g\, dy + h\, dz$ for $f,g,h \in C^{\infty}(U)$. For $\Omega^{2}(U)$, a form looks like $R\, dx \wedge dy + Q\, dx \wedge dz + P\, dy \wedge dz$ for $f,g,h \in C^{\infty}(U)$.
\end{example}

\textit{July 29th.}

\noindent The map $\Omega^{1}(U) \times \mf{X}(U) \to C^{\infty}(U)$, defined as $(\omega,X) \mapsto \omega(X)$, is $\R$-bilinear. In particular, for $\omega = \sum_{i=1}^{n} a_{i} dx^{i}$ and $X = \sum_{j=1}^{n} b^{j} \frac{\partial}{\partial x^{j}}$, we have
\begin{align}
    \omega(X) = \sum_{i=1}^{n} a_{i} dx^{i}\left( \sum_{j=1}^{n} b^{j} \frac{\partial}{\partial x^{j}} \right) = \sum_{i=1}^{n} a_{i} b^{i} \in C^{\infty}(U).
\end{align}
Extending this, we have the map $\Omega^{k}(U) \times \mf{X}(U)^{k} \to C^{\infty}(U)$, defined as $(\omega,X_{1},\ldots,X_{k}) \mapsto \omega(X_{1},\ldots,X_{k})$, which is $\R$-multilinear. In particular, for $\omega = \sum_{I} b_{I} dx^{I}$ and $X_{i} = \sum_{i=1}^{n} a_{i}^{j} \frac{\partial}{\partial x^{j}}$, we have
\begin{align}
    \omega(X_{1},\ldots,X_{k}) = \sum_{I} b_{I} a_{i_{1}}^{1} \cdots a_{i_{k}}^{k} \in C^{\infty}(U)
\end{align}
using the fact that $dx^{I}(\frac{\partial}{\partial x^{i_{1}}},\ldots,\frac{\partial}{\partial x^{i_{k}}}) = 1$ if the multi-index matches, and $0$ otherwise.

\subsection{Exterior Differentiation}

Here, $d_{k}$ shall denote the map $d_{k}:\Omega^{k}(U) \to \Omega^{k+1}(U)$, defined as the \eax{exterior derivative}. The subscript $k$ is often omitted, and we simply write $d:\Omega^{k}(U) \to \Omega^{k+1}(U)$. For $k = 0$, we have
\begin{align}
    d:C^{\infty}(U) \to \Omega^{1}(U), \quad f \mapsto df = \sum_{i=1}^{n} \frac{\partial f}{\partial x^{i}} dx^{i}.
\end{align}
Analogously, for a general $k$-form $\omega = \sum_{I} a_{I} dx^{I} \in \Omega^{k}(U)$, we have
\begin{align}
    d:\Omega^{k}(U) \to \Omega^{k+1}(U), \quad \omega \mapsto d\omega = \sum_{I} da_{I} \wedge dx^{I} = \sum_{I} \sum_{j=1}^{n} \frac{\partial a_{I}}{\partial x^{j}} dx^{j} \wedge dx^{I}.
\end{align}
All the $d$ maps are $\R$-linear, but not $C^{\infty}(U)$-linear. Also, for two forms $\omega \in \Omega^{k}(U)$ and $\eta \in \Omega^{l}(U)$, we have
\begin{align}
    d(\omega \wedge \eta) = d\omega \wedge \eta + (-1)^{k} \omega \wedge d\eta.
\end{align}
This is known as \eax{antiderivation}. Moreover, $d_{k+1} \circ d_{k} = d^{2} = 0$. Both these properties may be verified by the reader. Now suppose we have another map $D:\Omega^{k}(U) \to \Omega^{k+1}(U)$ which is $\R$-linear, satisfies antiderivation, and $D^{2} = 0$, and $Df = df$ for all $f \in C^{\infty}(U)$. Then $D = d$; the exterior derivative is unique. Thus, a chain of maps is formed:
\begin{align}
    0 \xrightarrow{d} C^{\infty}(U) \xrightarrow{d} \Omega^{1}(U) \xrightarrow{d} \Omega^{2}(U) \xrightarrow{d} \cdots \xrightarrow{d} \Omega^{n}(U) \xrightarrow{d} 0 \quad \text{ for open } U \subseteq \R^{n}.
\end{align}
An $\omega \in \Omega^{k}(U)$ is called an \eax{exact $k$-form} if $\omega = d\eta$ for some $\eta \in \Omega^{k-1}(U)$. It is called a \eax{closed $k$-form} if $d\omega = 0$. Note that every exact form is closed, but the converse is not true in general. We will refine the following later: quotienting the space of closed $k$-forms by the space of exact $k$-forms gives us the \eax{de Rham cohomology} of $U$, denoted by $\HdR^{k}(U)$.

\section{Manifolds}
\textit{July 31st.}

\begin{definition}
    A topological space $M$ is termed a \eax{topological manifold} of dimension $n$ if
    \begin{enumerate}
        \item $M$ is Hausdorff,
        \item $M$ is secound countable (that is, has a countable basis), and
        \item for all $p \in M$, there exists an open neighbourhood $U \subseteq M$ and a homemorphism $\varphi:U \to \varphi(U) \subseteq \R^{n}$, where $\varphi(U)$ is open. This makes a manifold \eax{locally euclidean}.
    \end{enumerate}
\end{definition}

The \eax{invariance of dimension} property states that if we have two open sets $U \subseteq \R^{n}$ and $V \subseteq \R^{m}$ such that $U$ is homeomorphic to $V$, then $m = n$. Thus, the dimension of a topological manifold is well-defined.

\begin{definition}
    A pair $(U,\varphi)$ is called a \eax{coordinate chart} on $M$ if $U \subseteq M$ is open and $\varphi:U \to \varphi(U) \subseteq \R^{n}$ is a homeomorphism. 
\end{definition}
A trivial example is $\R^{n}$ with one (coordinate) chart $(\R^{n},\id_{\R^{n}})$. In fact, if $U \subseteq \R^{n}$ is open, then $(U,\id_U)$ is a coordinate chart. Note that subspaces inherit the properties of Hausdorffness and second countability. We state some properties without proof; if $X$ is locally Euclidean but not $T_{2}$, then $X$ being a topological manifold implies that the number of connected (or path connected) components is at most countable.

\begin{definition}
    Two charts $(U,\varphi)$ and $(V,\psi)$ on a topological manifold $M$ are said to be \eax{$C^{\infty}$-compatible} if either
    \begin{enumerate}
        \item $U \cap V = \emptyset$, or
        \item $U \cap V \neq \emptyset$, and the transition maps $\psi \circ \varphi^{-1}:\varphi(U \cap V) \to \psi(U \cap V)$ and $\varphi \circ \psi^{-1}:\psi(U \cap V) \to \varphi(U \cap V)$ are $C^{\infty}$-maps between open subsets of $\R^{n}$.
    \end{enumerate}
    Similarly, $C^{k}$-compatibility may be defined for $k \geq 0$. A $C^{\omega}$-manifold is termed to be a real analytic manifold. In case of $k = 2m$, then $\R^{2m}$ is identified with $\C^{m}$, and the transition maps are required to be holomorphic. 
\end{definition}

\begin{definition}
    $M$ is a $C^{\infty}$-manifold or a \eax{smooth manifold} if $M$ is a topological manifold and there exists a family of charts $\{(U_{\alpha},\varphi_{\alpha})\}_{\alpha \in \Lambda}$ such that
    \begin{enumerate}
        \item $\{U_{\alpha}\}_{\alpha \in \Lambda}$ is an open cover of $M$, and
        \item $(U_{\alpha},\varphi_{\alpha})$ and $(U_{\beta},\varphi_{\beta})$ are $C^{\infty}$-compatible for all $\alpha,\beta \in \Lambda$.
    \end{enumerate}
    Similarly, $C^{k}$-manifolds, $C^{\omega}$-manifolds, and complex manifolds may be defined. The family of charts $\{(U_{\alpha},\varphi_{\alpha})\}_{\alpha \in \Lambda}$ is called an \eax{atlas} of $M$.
\end{definition}

The most basic non-trivial example is perhaps the circle $S^{1} = \{(x,y) \in \R^{2} \mid x^{2}+y^{2} = 1\}$, which is a $C^{\infty}$-manifold of dimension $1$. We can think of $\R^{2}$ as $\C$. Let $U = S^{1}\setminus\{1\} = \{e^{it} \mid 0 < t < 2\pi\}$ and $V = S^{1}\setminus\{-1\} = \{e^{it} \mid -\pi < t < \pi\}$. Then $(U,\varphi)$ and $(V,\psi)$ are charts, where $\varphi:U \to (0,2\pi) \subseteq \R$ is defined as $e^{it} \mapsto t$ and $\psi:V \to (-\pi,\pi) \subseteq \R$ is defined as $e^{it} \mapsto t$. Here, $U \cap V = \{e^{it} \mid 0 < t < \pi\} \cup \{e^{-t} \mid \pi < t < 2\pi,\, -\pi < t < 0\}$. The transition maps are $\psi \circ \varphi^{-1}:(0,\pi) \to (0,\pi)$ defined as $t \mapsto t$ and $\psi \circ \varphi^{-1}:(\pi,2\pi) \to (-\pi,0)$ defined as $t \mapsto t-2\pi$. Both these maps are smooth, and hence the charts are compatible. Thus, $S^{1}$ is a smooth manifold of dimension $1$.
\\

\textit{August 3rd.}
\begin{proposition}
    Let $\cU = \{(U_{\alpha},\varphi_{\alpha})\}_{\alpha \in \Lambda}$ be an atlas on $M$. Suppose $(V,\psi)$ and $(W,\sigma)$ are two charts on $M$ which are $C^{\infty}$-compatible with all charts in $\cU$. Then $(V,\psi)$ and $(W,\sigma)$ are $C^{\infty}$-compatible.
\end{proposition}
\begin{proof}
    Let $p \in V \cap W$. Then there exists $\alpha \in \Lambda$ such that $p \in U_{\alpha}$. Since $(V,\psi)$ and $(U_{\alpha},\varphi_{\alpha})$ are compatible, we have that $\psi \circ \varphi_{\alpha}^{-1}:\varphi_{\alpha}(U_{\alpha} \cap V) \to \psi(U_{\alpha} \cap V)$ is smooth. Similarly, since $(W,\sigma)$ and $(U_{\alpha},\varphi_{\alpha})$ are compatible, we have that $\sigma \circ \varphi_{\alpha}^{-1}:\varphi_{\alpha}(U_{\alpha} \cap W) \to \sigma(U_{\alpha} \cap W)$ is smooth. Now, we have
    \begin{align}
        (\sigma \circ \psi^{-1})|_{\psi(U_{\alpha} \cap V \cap W)} = (\sigma \circ \varphi_{\alpha}^{-1})|_{\varphi_{\alpha}(U_{\alpha} \cap V \cap W)} \circ (\varphi_{\alpha} \circ \psi^{-1})|_{\psi(U_{\alpha} \cap V \cap W)}.
    \end{align}
    The right-hand side is a composition of smooth maps, and hence is smooth. Thus, $\sigma \circ \psi^{-1}$ is smooth on an open neighbourhood of $p$. Since $p$ was arbitrary, we have that $\sigma \circ \psi^{-1}$ is smooth on $V\cap W$. Similarly, we can show that $\psi\circ\sigma^{-1}$ is smooth on $V\cap W$. Thus, $(V,\psi)$ and $(W,\sigma)$ are compatible.
\end{proof}

This leads to the following definition.

\begin{definition}
    A \eax{maximal atlas} $\cM$ on a smooth manifold $M$ is an atlas which is maximal with respect to inclusion; that is, if $\cM \subseteq \cU$ for some atlas $\cU$, then $\cM = \cU$. Equivalently, $\cM$ is a maximal atlas if it contains all charts which are compatible with all charts in $\cM$.
\end{definition}
Any atlas $\cU$ on $M$ is contained in a unique maximal atlas $\cM$. For a proof of this, let $\cM = \{(U,\varphi) \mid \text{ compatible with all charts in } \cU\}$. Then $\cM$ is an atlas. If $\cM \subseteq \cM'$, then any chart in $\cM'$ must be compatible with all charts in $\cM$; it must be compatible with all charts in $\cU$. By definition of $\cM$, such a chart must be in $\cM$, so $\cM = \cM'$. For uniqueness, if $\cU \subseteq \cM, \cM'$, then $\cM,\cM' \subseteq \cM \cup \cM'$ tells us $\cM = \cM'$ by maximality. The unique maximal atlas $\cM \supseteq \cU$ is also called a differentiable structure on $M$.

\begin{example}
    \begin{enumerate}
        \item Let $U \subseteq \R^{n}$, and $F:U \to \R^{m}$ be a $C^{\infty}$-map. The graph of $F$ is defined as $\Gamma(F) = \{(x,F(x)) \mid x \in U\} \subseteq \R^{n+m}$. At atlas on $\Gamma(F)$ is given by $\{(U,\varphi)\}$, where $\varphi:U \to \Gamma(F)$ is defined as $x \mapsto (x,F(x))$. The transition map is $\varphi^{-1}:\Gamma(F) \to U$ defined as $(x,F(x)) \mapsto x$. This is smooth, and hence $\Gamma(F)$ is a smooth manifold of dimension $n$.
        \item Let $M$ and $N$ be smooth manifolds of dimension $m$ and $n$, respectively. Then $M \times N$ is a smooth manifold of dimension $m+n$; if $\cU$ and $\cV$ are atlases on $M$ and $N$, respectively, then $\cU \times \cV = \{(U \times V,\varphi \times \psi) \mid (U,\varphi) \in \cU, (V,\psi) \in \cV\}$ is an atlas on $M \times N$. The transition maps are $(\varphi_{2} \times \psi_{2}) \circ (\varphi_{1} \times \psi_{1})^{-1} = (\varphi_{2} \circ \varphi_{1}^{-1}) \times (\psi_{2} \circ \psi_{1}^{-1})$, which is smooth since both $\varphi_{2} \circ \varphi_{1}^{-1}$ and $\psi_{2} \circ \psi_{1}^{-1}$ are smooth.
    \end{enumerate}
\end{example}



\subsection{Smooth Functions and Smooth Maps}

We will often omit `smooth' but still mean smooth in the context of manifolds.
\begin{definition}
    Let $M$ be a manifold of dimension $n$. A function $f:M \to \R$ is smooth at $p \in M$ if there exists a chart $(U,\varphi)$ of $M$ such that $f \circ \varphi^{-1}:\varphi(U) \to \R$ is smooth at $\varphi(p)$. A function $f:M \to \R$ is smooth if it is smooth at all $p \in M$.
\end{definition}
The definition of the smooth function $f:M \to \R$ is independent of the choice of chart $(U,\varphi)$, since if $(V,\psi)$ is another chart containing $p$, then $\psi \circ \varphi^{-1}:\varphi(U \cap V) \to \psi(U \cap V)$ is smooth, and hence $f \circ \psi^{-1} = (f \circ \varphi^{-1}) \circ (\varphi \circ \psi^{-1})$ is smooth at $\psi(p)$.

\begin{proposition}
    Let $M$ be a manifold of dimension $n$, and $f:M \to \R$ be a function. Then the following are equivalent:
    \begin{enumerate}
        \item $f$ is a smooth function on $M$.
        \item $M$ has an atlas $\cU = \{(U_{\alpha},\varphi_{\alpha})\}_{\alpha \in \Lambda}$ such that $f \circ \varphi_{\alpha}^{-1}:\varphi_{\alpha}(U_{\alpha}) \to \R$ is smooth for all $\alpha \in \Lambda$.
        \item For all charts $(U,\varphi)$ on $M$, the map $f \circ \varphi^{-1}:\varphi(U) \to \R$ is smooth.
    \end{enumerate}
\end{proposition}
\begin{proof}
    Note that 2.~implies 1. by definition, and 3.~implies 2. trivially. We show that 1.~implies 3. Let $(U,\varphi)$ be a chart on $M$, and $p \in U$; it is enough to prove that $f \circ \varphi^{-1}$ is smooth at $\varphi(p)$. Since $f$ is smooth, there exists a chart $(V,\psi)$ of $M$ such that $p \in V$ and $f \circ \psi^{-1}:\psi(V) \to \R$ is smooth at $\psi(p)$. Now, we have
    \begin{align}
        f \circ \varphi^{-1} = (f \circ \psi^{-1}) \circ (\psi \circ \varphi^{-1}).
    \end{align}
    The right-hand side is a composition of smooth maps, and hence is smooth at $\varphi(p)$. Thus, $f \circ \varphi^{-1}$ is smooth at $\varphi(p)$.
\end{proof}

A function $f:M \to \R$ is smooth on a chart $(U,\varphi)$ if and only if the pullback $\varphi^{\ast}f \defeq f \circ \varphi^{-1} :\varphi(U) \to \R$ is smooth. Also note that the definition of a smooth function $f$ does not assume that $f$ is continuous; however, one may show that the definition itself implies that $f$ is continuous. This is because if $f$ is smooth at $p \in M$, then there exists a chart $(U,\varphi)$ of $M$ such that $f \circ \varphi^{-1}:\varphi(U) \to \R$ is smooth at $\varphi(p)$, and hence continuous at $\varphi(p)$. Since $\varphi$ is a homeomorphism, we have that $f$ is continuous at $p$. Since $p$ was arbitrary, we have that $f$ is continuous on $M$.

\begin{definition}
    Let $M$ and $N$ be smooth manifolds of dimension $m$ and $n$, respectively. A continuous map $F:N \to M$ is smooth at $p \in N$ if there exist charts $(U,\varphi)$ on $N$, with $p \in U$, and $(V,\psi)$ on $M$, with $F(p) \in V$, such that $\psi \circ F \circ \varphi^{-1}:\varphi(U) \to \psi(V)$ is smooth at $\varphi(p)$. A continuous map $F:N \to M$ is smooth if it is smooth at all $p \in N$.
\end{definition}
The definition of a map $f$ being smooth at $p$ is independent of the choice of charts $(U,\varphi)$ and $(V,\psi)$, since if $(U',\varphi')$ and $(V',\psi')$ are other charts containing $p$ and $F(p)$, respectively, then $\psi' \circ \psi^{-1}:\psi(V \cap V') \to \psi'(V \cap V')$ and $\varphi' \circ \varphi^{-1}:\varphi(U \cap U') \to \varphi'(U \cap U')$ are smooth, and hence
\begin{align}
    \psi' \circ F \circ \varphi'^{-1} = (\psi' \circ \psi^{-1}) \circ (\psi \circ F \circ \varphi^{-1}) \circ (\varphi \circ \varphi'^{-1})
\end{align}

\begin{proposition}
    Let $F:N \to M$ be a continuous map between smooth manifolds. Then the following are equivalent:
    \begin{enumerate}
        \item $F$ is smooth.
        \item There exist atlases $\cU$ of $N$ and $\cV$ of $M$ such that for all $(U,\varphi) \in \cU$ and $(V,\psi) \in \cV$, the map $\psi \circ F \circ \varphi^{-1}:\varphi(U \cap F^{-1}(V)) \to \psi(V)$ is smooth.
        \item For all charts $(U,\varphi)$ on $N$ and $(V,\psi)$ on $M$, the map $\psi \circ F \circ \varphi^{-1}:\varphi(U \cap F^{-1}(V)) \to \psi(V)$ is smooth.
    \end{enumerate}
\end{proposition}

\begin{proof}
    Again, note that 2.~implies 1. by definition, and 3.~implies 2. trivially. A similar argument as in the previous proposition shows that 1.~implies 3.
\end{proof}

\begin{proposition}
    Let $F:N \to M$ and $G:M \to P$ be smooth maps between smooth manifolds. Then the composition $G \circ F:N \to P$ is smooth.
\end{proposition}
\begin{proof}
    Let $(U,\varphi)$, $(V,\psi)$, and $(W,\sigma)$ be charts on $N$, $M$, and $P$, respectively, such that $p \in U$, $F(p) \in V$, and $G(F(p)) \in W$. Then we have 
    \begin{align}
        \sigma \circ (G \circ F) \circ \varphi^{-1} = (\sigma \circ G \circ \psi^{-1}) \circ (\psi \circ F \circ \varphi^{-1}).
    \end{align}
    The right-hand side is a composition of smooth maps, and hence is smooth at $\varphi(p)$. Thus, $G \circ F$ is smooth at $p$.
\end{proof}

Finally, a smooth map $F:N \to M$ is termed a \eax{diffeomorphism} if $F$ is a bijection, and $F$ and $F^{-1}$ are smooth. Some properties of diffeomorphisms may be inferred.

\begin{proposition}
    The following hold true.
    \begin{enumerate}
        \item Let $(U,\varphi)$ be a chart on a smooth manifold $M$ of dimension $n$. Then $\varphi:U \to \varphi(U)$ is a diffeomorphism.
        \item Suppose $U \subseteq M$ is open, and that $F:U \to F(U)$ is a diffeomorphism. Then $(U,F)$ is a chart in the maximal atlas of $M$.
    \end{enumerate}
\end{proposition}
\begin{proof}
    \begin{enumerate}
        \item Note that $U$ is a manifold since it is open in $M$, and $\varphi(U)$ is a manifold since it is open in $\R^{n}$. The map $\varphi:U \to \varphi(U)$ is a homeomorphism by definition of a chart, so it is bijective. Note that $\{(U,\varphi)\}$ is an atlas on $U$, and $\{(\varphi(U),\id_{\varphi(U)})\}$ is an atlas on $\varphi(U)$. Then $\varphi$ is smooth since $\id_{\varphi(U)} \circ \varphi \circ \varphi^{-1} = \id_{\varphi(U)}$ is smooth, and $\varphi^{-1}$ is smooth since $\varphi \circ \varphi^{-1} \circ \id_{\varphi(U)} = \id_{\varphi(U)}$ is smooth. Thus, $\varphi$ is a diffeomorphism.
        
        \item Since $U$ is open in $M$, it is a smooth manifold. Because $F:U\to F(U)$ is a diffeomorphism, it is in particular a homeomorphism. Hence $F$ is bijective, and its image $F(U)$ is an open subset of $\R^{n}$ (as required for a chart). Let $(V,\psi)$ be any chart on $M$ with $U\cap V\neq\varnothing$. Then the transition maps
        \begin{align}
            \psi\circ F^{-1}:F(U\cap V)\to \psi(V) \quad \text{and} \quad F\circ \psi^{-1}:\psi(U\cap V)\to F(U\cap V)
        \end{align}
        are smooth, since they are compositions of smooth maps. Therefore $(U,F)$ is compatible with every chart in the maximal atlas of $M$, and so $(U,F)$ is a chart in that maximal atlas.
    \end{enumerate}
\end{proof}

\begin{proposition}
    Let $N$ be a smooth manifold of dimension $n$, and $F:N \to \R^{m}$ be a continuous map. Then the following are equivalent.
    \begin{enumerate}
        \item $F$ is smooth.
        \item $N$ has an atlas $\cU$ such that for all $(U,\varphi) \in \cU$, the map $F \circ \varphi^{-1}:\varphi(U) \to \R^{m}$ is smooth.
        \item For all charts $(U,\varphi)$ on $N$, the map $F \circ \varphi^{-1}:\varphi(U) \to \R^{m}$ is smooth.
    \end{enumerate}
\end{proposition}
\begin{proof}
    3.~implies 2.~is obvious. For 2.~implies 1., note that there exist atlases $\cU$ and $\cV$ of $N$ and $M$ respectively such that $(U,\varphi) \in \cU$ and $(V,\psi) \in \cV$ satisfy the smoothness condition. Since $\R^{m}$ has a global chart $(\R^{m},\id_{\R^{m}})$, we can take $\cV = \{(\R^{m},\id_{\R^{m}})\}$. Then for any $(U,\varphi) \in \cU$, we have $\id_{\R^{m}} \circ F \circ \varphi^{-1} = F \circ \varphi^{-1}$ is smooth, which implies that $F$ is smooth. For 1.~implies 3., $F$ is smooth and $\varphi$ is a diffeomorphism; thus, $\varphi^{-1}$ is smooth, and $F \circ \varphi^{-1}$ is smooth.
\end{proof}

\begin{proposition}
    Let $F:N \to M$ be a continuous map between smooth manifolds. Then the following are equivalent.
    \begin{enumerate}
        \item $F$ is smooth.
        \item $M$ has an atlas $\cV$ such that for all $(V,\psi) \in \cV$, the map $\psi \circ F:F^{-1}(V) \to \psi(V)$ is smooth.
        \item For any chart $(V,\psi)$ of $M$, the map $\psi \circ F:F^{-1}(V) \to \psi(V)$ is smooth.
    \end{enumerate}
\end{proposition}
The proof follows a similar structure to the one above.

\begin{proposition}
    Let $F:N \to \R^{m}$ be a continuous map. Then the following are equivalent.
    \begin{enumerate}
        \item $F$ is smooth.
        \item $F^{1},\ldots,F^{m}:N \to \R$ are smooth, where $F = (F^{1},\ldots,F^{m})$.
    \end{enumerate}
\end{proposition}

\subsubsection{Examples}

\textit{August 5th.}

We look at some examples of smooth manifolds and smooth maps.

\begin{example}
    For $S^{n} \subseteq \R^{n+1}$, the \eax{$n$-sphere}, define $U_{i}^{+} = \{(x_{1},\ldots,x_{n+1}) \in S^{n} \mid x_{i} > 0\}$ and $U_{i}^{-} = \{(x_{1},\ldots,x_{n+1}) \in S^{n} \mid x_{i} < 0\}$ for $i = 1,\ldots,n+1$. Then $\{(U_{i}^{+},\varphi_{i}^{+}), (U_{i}^{-},\varphi_{i}^{-})\}_{i=1}^{n+1}$ is an atlas on $S^{n}$, where $\varphi_{i}^{+}:U_{i}^{+} \to \R^{n}$ is defined as $(x_{1},\ldots,x_{n+1}) \mapsto (x_{1},\ldots,x_{i-1},x_{i+1},\ldots,x_{n+1})$, and $\varphi_{i}^{-}:U_{i}^{-} \to \R^{n}$ is defined similarly. The transition maps are smooth, and hence $S^{n}$ is a smooth manifold of dimension $n$.
\end{example}

\begin{example}
    Recall $\R P^{n}$, the \eax{real projective space} of dimension $n$, defined as the quotient $\R^{n+1}\setminus\{0\}/\sim$, where $(x_{1},\ldots,x_{n+1}) \sim (y_{1},\ldots,y_{n+1})$ if there exists $\lambda \in \R\setminus\{0\}$ such that $(x_{0},\ldots,x_{n}) = (\lambda y_{0},\ldots,\lambda y_{n})$. Let $(a_{1}:a_{2}:\ldots:a_{n+1})$ denote the equivalence class of $(a_{1},\ldots,a_{n+1}) \in \R^{n+1}\setminus\{0\}$. This is a smooth manifold of dimension $n$; an atlas is given by $\{(U_{i},\varphi_{i})\}_{i=1}^{n+1}$, where $U_{i} = \{(x_{1}:\ldots:x_{n+1}) \in \R P^{n} \mid x_{i} \neq 0\}$ and $\varphi_{i}:U_{i} \to \R^{n}$ is defined as
    \begin{align}
        \varphi_{i}(x_{1}:\ldots:x_{n+1}) = \left( \frac{x_{1}}{x_{i}},\ldots,\frac{x_{i-1}}{x_{i}},\frac{x_{i+1}}{x_{i}},\ldots,\frac{x_{n+1}}{x_{i}} \right).
    \end{align}
    Verify that the transition maps $\varphi_{j} \circ \varphi_{i}^{-1}:\varphi_{i}(U_{i} \cap U_{j}) \to \varphi_{j}(U_{i} \cap U_{j})$ are smooth for all $i,j = 1,\ldots,n+1$. We can also provide a smooth map $F:S^{n} \to \R P^{n}$ defined as $(x_{1},\ldots,x_{n+1}) \mapsto (x_{1}:\ldots:x_{n+1})$. This is a smooth map, and it is a $2$-to-$1$ covering map.
\end{example}

\begin{example}
    Similar to the above, we can define $\Prob_{k}^{n}$ to be the projective space of dimension $n$ over the field $k$; that is, $\Prob_{k}^{n} = k^{n+1}\setminus\{0\}/\sim$, where $(x_{1},\ldots,x_{n+1}) \sim (y_{1},\ldots,y_{n+1})$ if there exists $\lambda \in k\setminus\{0\}$ such that $(x_{1},\ldots,x_{n+1}) = (\lambda y_{1},\ldots,\lambda y_{n+1})$. In particular, if $k = \C$, then $\Prob_{\C}^{n}$ consists of complex 1-dimensional subspaces of $\C^{n+1}$. This is a smooth manifold of dimension $2n$. $\R P^{n}$ may be, thus, interpreted as the space of real 1-dimensional subspaces of $\R^{n+1}$. Building upon this, the \eax{Grassmannian space} $\Gr(n,k)$ is defined as the space of $k$-dimensional subspaces of $\R^{n}$. This is a smooth manifold of dimension $k(n-k)$. The Grassmannian space $\Gr(n+1,1)$ is precisely $\R P^{n}$, and $\Gr(n,0)$ is a single point. $\Gr(n,n-1)$ is the space of hyperplanes in $\R^{n}$. Again, the Grassmannian construction may be generalized to any field.
\end{example}

\subsection{Quotient Manifolds}
\textit{August 10th.}

Recall that for an equivalence relation $\sim$ on a set $S$, the quotient set $S/\!\sim$ is defined as the set of equivalence classes $[s] = \{t \in S \mid t \sim s\} \in S/\!\sim$. If $S$ happens to be a topological space, then the projection map $\pi:S \to S/\!\sim$ defined as $s \mapsto [s]$ is continuous if we endow $S/\!\sim$ with the quotient topology; that is, a subset $U \subseteq S/\!\sim$ is open if and only if $\pi^{-1}(U) \subseteq S$ is open. This is the finest topology on $S/\!\sim$ which makes $\pi$ continuous. Also recall the universal property; that is, if $f:S \to X$ is a set map between topological spaces, and $\pi:S \to S/\!\sim$ is the projection map, then there exists a unique map $\tilde{f}:S/\!\sim \to X$ such that $f = \tilde{f} \circ \pi$ if and only if $f$ is constant on equivalence classes. In this case, $\tilde{f}$ is continuous if and only if $f$ is continuous.

The equivalence relation $\sim$ on $S$ is called open if the projection map $\pi:S \to S/\!\sim$ is open; that is, if $U \subseteq S$ is open, then $\pi(U) \subseteq S/\!\sim$ is open. Usually, $S/\!\sim$ is not $T_{2}$ is general. The following is a sufficient condition for $S/\!\sim$ to be $T_{2}$, or Hausdorff.

\begin{proposition}
    Let $\sim$ be an open equivalence relation on a topological space $S$. Then $S/\!\sim$ is $T_{2}$ if and only if $\{(x,y) \in S \times S \mid x \sim y\}$ is closed in $S \times S$.
\end{proposition}

\begin{proof}
    Suppose the set $R = \{(x,y) \in S \times S \mid x \sim y\}$ is closed in $S \times S$. That is, for all $(x,y) \in S \times S$ such that $x \not\sim y$, there exist open neighbourhoods $U$ of $x$ and $V$ of $y$ such that $U \times V \subseteq S \times S \setminus R$. Since $\sim$ is open, we have that $\pi(U)$ and $\pi(V)$ are open neighbourhoods of $[x]$ and $[y]$, respectively. We claim that $\pi(U) \cap \pi(V) = \emptyset$. Suppose not; then there exists $[z] \in \pi(U) \cap \pi(V)$. Then there exist $u \in U$ and $v \in V$ such that $[z] = [u] = [v]$, which implies that $u \sim v$. But this contradicts the fact that $(u,v) \in U \times V \subseteq S \times S \setminus R$. Thus, $\pi(U) \cap \pi(V) = \emptyset$, and hence $S/\!\sim$ is Hausdorff.

    On the other hand, suppose $S/\!\sim$ is Hausdorff. Let $(x,y) \in S \times S$ such that $x \not\sim y$. Then there exist open neighbourhoods $U$ of $[x]$ and $V$ of $[y]$ in $S/\!\sim$ such that $U \cap V = \emptyset$. Since $\pi:S \to S/\!\sim$ is continuous, we have that $\pi^{-1}(U)$ and $\pi^{-1}(V)$ are open neighbourhoods of $x$ and $y$, respectively. We claim that $\pi^{-1}(U) \times \pi^{-1}(V) \subseteq S \times S \setminus R$. Suppose not; then there exists $(u,v) \in \pi^{-1}(U) \times \pi^{-1}(V)$ such that $u \sim v$. Then we have $[u] = [v]$, which implies that $[u] \in U \cap V$, contradicting the fact that $U \cap V = \emptyset$. Thus, $\pi^{-1}(U) \times \pi^{-1}(V) \subseteq S \times S \setminus R$, and hence $R$ is closed in $S \times S$.
\end{proof}


\begin{example}
    Let $\R^{n+1} \setminus \{(0,\ldots,0)\}$ be the punctured space, and define an equivalence relation $\sim$ on it as $(x_{1},\ldots,x_{n+1}) \sim (y_{1},\ldots,y_{n+1})$ if there exists $\lambda \in \R\setminus\{0\}$ such that $(x_{1},\ldots,x_{n+1}) = (\lambda y_{1},\ldots,\lambda y_{n+1})$. Then $\R^{n+1} \setminus \{(0,\ldots,0)\}/\!\sim$ is precisely $\R P^{n}$. The projection map $\pi:\R^{n+1} \setminus \{(0,\ldots,0)\} \to \R P^{n}$ is open; let $U \subseteq \R^{n}\setminus\{(0,\ldots,0)\}$ be open. Then $\pi^{-1}(\pi(U)) = \bigcup_{\lambda \in \R\setminus\{0\}} \lambda U$ is open, and hence $\pi(U)$ is open in $\R P^{n}$. Also, the set $R = \{(x,y) \in (\R^{n+1}\setminus\{0\}) \times (\R^{n+1}\setminus\{0\}) \mid x \sim y\}$ is closed in $(\R^{n+1}\setminus\{0\}) \times (\R^{n+1}\setminus\{0\})$, since if $y_{i}=\lambda x_{i}$ for some $\lambda \in \R\setminus\{0\}$, then $x_{i}y_{j}-x_{j}y_{i}=0$ for all $i,j=1,\ldots,n+1$. Thus, $R$ is the zero set of the continuous functions $f_{ij}:(\R^{n+1}\setminus\{0\}) \times (\R^{n+1}\setminus\{0\}) \to \R$ defined as $(x,y) \mapsto x_{i}y_{j}-x_{j}y_{i}$ for all $i,j=1,\ldots,n+1$. Hence, $R$ is closed in $(\R^{n+1}\setminus\{0\}) \times (\R^{n+1}\setminus\{0\})$. By the previous proposition, we have that $\R P^{n}$ is Hausdorff.
\end{example}

\begin{proposition}
    Let $\sim$ be an open equivalence relation on a topological space $S$. Then $S$ is second countable if and only if $S/\!\sim$ is second countable.
\end{proposition}
\begin{proof}
    
    Suppose $S$ is second countable, and let $\cB$ be a countable basis for $S$. Then $\cB' = \{\pi(B) \mid B \in \cB\}$ is a countable collection of open sets in $S/\!\sim$. We claim that $\cB'$ is a basis for $S/\!\sim$. Let $U \subseteq S/\!\sim$ be open, and let $[x] \in U$. Then $\pi^{-1}(U) \subseteq S$ is open, and hence there exists $B \in \cB$ such that $x \in B \subseteq \pi^{-1}(U)$. Then we have $[x] \in \pi(B) \subseteq U$, and hence $\cB'$ is a basis for $S/\!\sim$. Thus, $S/\!\sim$ is second countable.

    On the other hand, suppose $S/\!\sim$ is second countable, and let $\cB'$ be a countable basis for $S/\!\sim$. Then $\cB = \{\pi^{-1}(B') \mid B' \in \cB'\}$ is a countable collection of open sets in $S$. We claim that $\cB$ is a basis for $S$. Let $U \subseteq S$ be open, and let $x \in U$. Then $\pi(U) \subseteq S/\!\sim$ is open, and hence there exists $B' \in \cB'$ such that $[x] \in B' \subseteq \pi(U)$. Then we have $x \in \pi^{-1}(B') \subseteq U$, and hence $\cB$ is a basis for $S$. Thus, $S$ is second countable.
\end{proof}

\begin{example}
    Consider the disc $D^{n} = \{(x_{1},\ldots,x_{n}) \in \R^{n} \mid x_{1}^{2}+\ldots+x_{n}^{2} \leq 1\}$, and the hemisphere $H^{n} = \{(x_{1},\ldots,x_{n+1}) \in S^{n} \mid x_{n+1} \geq 0\}$. On $H^{n}$, define an equivalence relation as $(x_{1},\ldots,x_{n},x_{n+1}) \sim (y_{1},\ldots,y_{n},y_{n+1})$ if either
    \begin{enumerate}
        \item $x_{n+1} = y_{n+1} = 0$ and $(y_{1},\ldots,y_{n},0) = (-x_{1},\ldots,-x_{n},0)$, or
        \item $(x_{1},\ldots,x_{n},x_{n+1}) = (y_{1},\ldots,y_{n},y_{n+1})$.
    \end{enumerate}
    On $D^{n}$, define an equivalence relation simply as $x \sim y$ if $x = y$ or $x = -y$. Then, one can show that $\R P^{n} \cong S^{n}/\!\sim \; \cong D^{n}/\!\sim \; \cong H^{n}/\!\sim$ as smooth manifolds, where $\sim$ on $S^{n}$ is the antipodal equivalence relation.
\end{example}

\section{Differential Structure}
Let $M$ be a smooth manifold of dimension $n$. For a point $p \in M$, let $(U,\varphi)$ be a chart of $M$ such that $p \in U$. The functions $r^{1},\ldots,r^{n}:\R^{n} \to \R$ defined as $r^{i}(x_{1},\ldots,x_{n}) = x_{i}$ are smooth, and hence, we can interpret $(U,\varphi)$ as $(U,x^{1},\ldots,x^{n})$, a chart with smooth coordinate functions $x^{i} = r^{i} \circ \varphi:U \to \R$ for $i = 1,\ldots,n$. Then, given a map $f \in C^{\infty}(U)$, we define the partial derivative of $f$ with respect to the coordinate function $x^{i}$ at $p$ as
\begin{align}
    \frac{\partial f}{\partial x^{i}}(p) \defeq \frac{\partial (f \circ \varphi^{-1})}{\partial r^{i}}(\varphi(p)).
\end{align}
In fact, given any $p$ in this particular $U$, we have
\begin{align}
    (\varphi^{-1})^{\ast}\left(\frac{\partial f}{\partial x^{i}} \right) = \frac{\partial f}{\partial x^{i}} \circ \varphi^{-1} = \frac{\partial (f \circ \varphi^{-1})}{\partial r^{i}} \in C^{\infty}(\varphi(U)).
\end{align}
This implies that $\frac{\partial f}{\partial x^{i}} \in C^{\infty}(U)$, and hence the partial derivative of $f$ with respect to $x^{i}$ is smooth on $U$. Moreover, one can show that $\frac{\partial x^{i}}{\partial x^{j}}(p) = \frac{\partial r^{i}}{\partial r^{j}}(\varphi(p)) = \delta_{j}^{i}$.

Let $F:N \to M$ be a smooth map between smooth manifolds of dimension $n$ and $m$, respectively. Let $(U,\varphi)$ be a chart on $N$ with $p \in U$, and $(V,\psi)$ be a chart on $M$ with $F(p) \in V$. Then we can interpret $(U,\varphi)$ as $(U,x^{1},\ldots,x^{n})$, and $(V,\psi)$ as $(V,y^{1},\ldots,y^{m})$, where $x^{i} = r^{i} \circ \varphi:U \to \R$ and $y^{j} = r^{j} \circ \psi:V \to \R$ are smooth coordinate functions. Then, we have $F^{i} = y^{i} \circ F = r^{i} \circ \psi \circ F \in C^{\infty}(U)$ for $i = 1,\ldots,m$. The \eax{Jacobian matrix} of $F$ at $p$, with respect to the charts $(U,\varphi)$ and $(V,\psi)$, is defined as the $m \times n$ matrix
\begin{align}
    J_{F}(p) = \left( \frac{\partial F^{i}}{\partial x^{j}}(p) \right)_{1 \leq i \leq m, 1 \leq j \leq n}.
\end{align}
Each entry here can be written as
\begin{align}
    \frac{\partial F^{i}}{\partial x^{j}}(p) = \frac{\partial (\psi \circ F \circ \varphi^{-1})^{i}}{\partial r^{j}}(\varphi(p)) = \frac{\partial (r^{i}\psi F \varphi^{-1})}{\partial r^{j}}(\varphi(p)) = \frac{\partial(y^{i} F \varphi^{-1})}{\partial r^{j}}(\varphi(p)) = \frac{\partial(F^{i} \circ \varphi^{-1})}{\partial r^{j}}(\varphi(p)).
\end{align}
Recall the inverse function theorem, which states that if $F:U\subseteq\R^{n} \to \R^{n}$ is smooth, and the Jacobian matrix $J_{F}(p)$ is invertible at $p \in U$, then there exists an open neighbourhood $V$ of $F(p)$ such that $F:U \to V$ is a diffeomorphism. This motivates the following definition. Thus, if $m = n$ and $J_{F}(p)$ is invertible, then $F$ is a local diffeomorphism at $p$; $\psi circ F \circ \varphi^{-1}$ is a local diffeomorphism at $\varphi(p)$ if and only if $F$ is a local diffeomorphism at $p$.

Let $M$ be a smooth manifold of dimension $n$, and let $p \in M$. Suppose $F^{1},F^{2},\ldots,F^{n}$ are smooth functions on some open neighbourhood $U \subseteq M$ of $p$. Then consider $F = (F^{1},\ldots,F^{n}):U \to \R^{n}$, which is a smooth map. Then $\left(\frac{\partial F^{i}}{\partial x^{j}}(p)\right)_{n \times n}$ is invertible if and only if $F$ is a local diffeomorphism at $p$ if and only if $F:W \to F(W) \subseteq \R^{n}$ is a diffeomorphism for some open neighbourhood $W \subseteq U$ of $p$. Then, on some open set around $p$, we can use $F^{1},\ldots,F^{n}$ as local coordinates. This motivates the following definition.


\begin{definition}
    Let $M$ be a smooth manifold, with $p \in M$. Recall $C_{p}^{\infty}(M) \equiv C_{p}^{\infty}$, the local ring of germs of smooth functions at $p$, which is also an $\R$-algebra. Also define $T_{p}M$ to be the $\R$-vector space of derivations at $p$; that is, $T_{p}M = \{D:C_{p}^{\infty} \to \R \mid D \text{ is linear and satisfies the Leibniz rule}\}$. This $T_{p}M$ is called the \eax{tangent space} of $M$ at $p$.
\end{definition}

\noindent \textit{August 12th.}

We leave it as an exercise to check that $\frac{\partial}{\partial x^{i}}|_{p}$ is a derivation from $C_{p}^{\infty}$ to $\R$, and hence $\frac{\partial}{\partial x^{i}}|_{p} \in T_{p}M$ for all valid $i$. Given a smooth map $F:N \to M$ between smooth manifolds, we can define the map $F_{\ast,p}:T_{p}N \to T_{F(p)}M$ as $D \mapsto F_{\ast,p}(D)$, where $F_{\ast,p}(D)(f) = D(f \circ F)$ for all $f \in C_{F(p)}^{\infty}$. This is the pushforward of $f$ at $p$. Moreover, if $G:M \to P$ is another smooth map, then $(G \circ F)_{\ast,p} = G_{\ast,F(p)} \circ F_{\ast,p}$. The map $F_{\ast,p}$ is linear, and hence a linear transformation between the tangent spaces $T_{p}N$ and $T_{F(p)}M$.

If $F:N \to M$ were a diffeomorphism, then there exists a diffeomorphism $G:M \to N$ such that $G \circ F = \id_{N}$ and $F \circ G = \id_{M}$. Now look at the chain
\begin{align}
    T_{p}N \xrightarrow{F_{\ast,p}} T_{F(p)}M \xrightarrow{G_{\ast,F(p)}} T_{G(F(p))}N = T_{p}N.
\end{align}
Here, we have $(G \circ F)_{\ast,p} = G_{\ast,F(p)} \circ F_{\ast,p} = (\id_{N})_{\ast,p} = \id_{T_{p}N}$, and similarly $F_{\ast,p} \circ G_{\ast,F(p)} = \id_{T_{F(p)}M}$. Thus, $F_{\ast,p}$ is a linear isomorphism between the tangent spaces $T_{p}N$ and $T_{F(p)}M$.
\\

Let $p \in M$, a smooth manifold, and let $(U,\varphi)$ be a chart containing $p$. Since $C_{p}^{\infty}(M) \equiv C_{p}^{\infty}(U)$, the diffeomorphism $\varphi:U \to \varphi(U)$ induces an isomorphism $\varphi_{\ast,p}:T_{p}M \to T_{\varphi(p)}\R^{n}$. We know that the tangent space on the right has the basis $\left\{\frac{\partial}{\partial r^{1}}|_{\varphi(p)},\ldots,\frac{\partial}{\partial r^{n}}|_{\varphi(p)}\right\}$. Applying this map to a basis element, we have
\begin{align}
    \varphi_{\ast,p}\left(\frac{\partial}{\partial x^{i}}\Big|_{p}\right)(f) = \frac{\partial}{\partial x^{i}}\Big|_{p}(f \circ \varphi) = \frac{\partial}{\partial r^{i}}\Big|_{\varphi(p)}(f \circ \varphi \circ \varphi^{-1}) = \frac{\partial}{\partial r^{i}}\Big|_{\varphi(p)}(f) \quad \text{for all } f \in C_{p}^{\infty}(M).
\end{align}
Hence, $\varphi_{\ast,p}\left(\frac{\partial}{\partial x^{i}}|_{p}\right) = \frac{\partial}{\partial r^{i}}|_{\varphi(p)}$, and so $\left\{\frac{\partial}{\partial x^{1}}|_{p},\ldots,\frac{\partial}{\partial x^{n}}|_{p}\right\}$ is a basis for $T_{p}M$. Thus, we have shown that the tangent space $T_{p}M$ has dimension $n$. This also shows that if $F:N \to M$ is a smooth map between smooth manifolds of dimension $n$ and $m$ that also happens to be a diffeomorphism, then $n = m$. Moreover, given a map $F:N \to M$ between smooth manifolds, and given a chart $(U,\varphi)$ of $p \in N$ and a chart $(V,\psi)$ of $F(p) \in M$, we have
\begin{align}
    F_{\ast,p}:T_{p}N \to T_{F(p)}M, \quad F_{\ast,p}\left(\frac{\partial}{\partial x^{i}}\Big|_{p}\right) = \sum_{j=1}^{m} b_{j}^{i} \frac{\partial}{\partial y^{j}}\Big|_{F(p)}, \quad b_{j}^{i} = \frac{\partial F^{j}}{\partial x^{i}}(p) \quad \text{for } i = 1,\ldots,n, j = 1,\ldots,m.
\end{align}
That is, the matrix of $F_{\ast,p}$ with respect to the bases $\left\{\frac{\partial}{\partial x^{1}}|_{p},\ldots,\frac{\partial}{\partial x^{n}}|_{p}\right\}$ and $\left\{\frac{\partial}{\partial y^{1}}|_{F(p)},\ldots,\frac{\partial}{\partial y^{m}}|_{F(p)}\right\}$ is precisely the $m \times n$ real \eax{Jacobian matrix} $(dF)_{p} \defeq J_{F}(p) = \left( \frac{\partial F^{i}}{\partial x^{j}}(p) \right)_{1 \leq i \leq m, \, 1 \leq j \leq n}$.
\\

\noindent \textit{August 17th.}

Suppose $0 \in (a,b) \subseteq \R$, and let $c:(a,b) \to M$ be a smooth map to a smooth manifold $M$ of dimension $n$ such that $c(0) = p \in M$. $c$ is then called a \eax{smooth curve} in $M$. Let $(U,\varphi) = (U,x^{1},\ldots,x^{n})$ be a chart on $M$ with $p \in U$. Then $c^{1}(0) \in T_{p}M$ is defined as $c'(0) = c_{\ast,0}\left(\frac{d}{dt}\Big|_{0}\right)$, where $\frac{d}{dt}|_{0} \in T_{0}\R$ is the standard basis element. Then $c'(0)$ is a tangent vector at $p$, and we have
\begin{align}
    \R \cong T_{0}((a,b)) \xrightarrow{c_{\ast,0}} T_{p}M \xrightarrow{\varphi_{\ast,p}} T_{\varphi(p)}\R^{n} \cong \R^{n}, \quad c'(0) \mapsto \varphi_{\ast,p}(c'(0)) = (\varphi \circ c)_{\ast,0}\left(\frac{d}{dt}\Big|_{0}\right) = \sum_{i=1}^{n} a^{i} \frac{\partial}{\partial x^{i}}\Big|_{p}
\end{align}
since $\{\frac{\partial}{\partial x^{i}} |_{p}\}_{i=1}^{n}$ is a basis for $T_{p}M$. Utilising $c'(0) = c_{\ast,0}\left(\frac{d}{dt}|_{0}\right)$, we have
\begin{align}
    a^{i} = \varphi_{\ast,p}(c'(0))(x^{i}) = c_{\ast,0}\left(\frac{d}{dt}\Big|_{0}\right)(x^{i} \circ \varphi^{-1}) = \frac{d}{dt}\Big|_{0}(x^{i} \circ \varphi^{-1} \circ c) = \frac{d}{dt}\Big|_{0}(x^{i} \circ c) = \frac{dc^{i}}{dt}(0).
\end{align}
Thus, a curve $c$ can be thought of as a tangent vector at $p$, and the tangent vector is given by the derivative of the curve at $0$. In particular, let us interpret $c:\R \to \R^{n}$ as a smooth curve in $\R^{n}$; we have $c(t) = (c_{1}(t),\ldots,c_{n}(t))$, and hence $c'(0) = (\frac{dc_{1}}{dt}(0),\ldots,\frac{dc_{n}}{dt}(0)) \in \R^{n}$. This is consistent with the above definition of $c'(0)$ as a tangent vector at $p$. From here, we wish to obtain a smooth map $c:(a,b) \to M$ such that $c(0) = p$ and $c'(0) = X_{p} = \sum_{i=1}^{n} a^{i} \frac{\partial}{\partial x^{i}}|_{p}$ for some $X_{p} \in T_{p}M$. Let $(U,\varphi) = (U,x^{1},\ldots,x^{n})$ be a chart on $M$ with $p \in U$. In fact, choose a chart $\varphi$ such that $\varphi(p) = 0 \in \R^{n}$ (translations are diffeomorphisms). Then we can define a smooth curve $c:(-\epsilon,\epsilon) \to M$ for some $\epsilon > 0$ such that $c(0) = p$ and $c'(0) = X_{p}$ as
\begin{align}
    c(t) = \varphi^{-1}(ta^{1},\ldots,ta^{n}) \quad \text{for } t \in (-\epsilon,\epsilon).
\end{align}
Then, $c'(0) = c_{\ast,0}\left(\frac{d}{dt}|_{0}\right) = \sum_{i=1}^{n} a^{i} \frac{\partial}{\partial x^{i}}|_{p} = X_{p}$, as desired. 

Coming back to the smooth map $F:N \to M$ between smooth manifolds, the rank of $F$ at $p \in N$ is defined as the rank of the linear map $F_{\ast,p}:T_{p}N \to T_{F(p)}M$, which is also the rank of the matrix $(dF)_{p} = J_{F}(p) = \left( \frac{\partial F^{i}}{\partial x^{j}}(p) \right)_{1 \leq i \leq m, 1 \leq j \leq n}$. The map $F:N \to M$ is called an \eax{immersion} at $p \in N$ if $F_{\ast,p}:T_{p}N \to T_{F(p)}M$ is injective (that is, $\rank(F_{\ast,p}) = \dim(T_{p}N) = n$), and $F$ is called an immersion if it is an immersion at all points $p \in N$. Similarly, $F$ is called a \eax{submersion} at $p \in N$ if $F_{\ast,p}:T_{p}N \to T_{F(p)}M$ is surjective (that is, $\rank(F_{\ast,p}) = \dim(T_{F(p)}M) = m$), and $F$ is called a submersion if it is a submersion at all points $p \in N$. If $F:N \to M$ is a smooth map between smooth manifolds of dimension $n$ and $m$, respectively, then $\rank(F_{\ast,p}) \leq \min\{n,m\}$ for all $p \in N$. Thus, if $n < m$, then $F$ cannot be a submersion, and if $n > m$, then $F$ cannot be an immersion.

\section{Regular Maps and Regular Submanifolds}

\begin{definition}
    Let $F:N \to M$ be a smooth map between smooth manifolds. $p \in N$ is called a \eax{critical point} of $F$ if $F_{\ast,p}:T_{p}N \to T_{F(p)}M$ is not surjective; that is, $F$ is not a submersion at $p$. $p$ is called a \eax{regular point} of $F$ if $F_{\ast,p}:T_{p}N \to T_{F(p)}M$ is surjective; that is, $F$ is a submersion at $p$. Similarly, $q \in M$ is called a \eax{critical value} of $F$ if there exists a critical point $p \in N$ such that $F(p) = q$. $q \in M$ is called a \eax{regular value} of $F$ if it is not a critical value; that is, for all $p \in N$ such that $F(p) = q$, we have that $p$ is a regular point of $F$. That is, either $F^{-1}(m) = \emptyset$, or $F_{\ast,p}:T_{p}N \to T_{F(p)}M$ is surjective for all $p \in F^{-1}(m)$.
\end{definition}

From here, we can define something called a regular submanifold, which is useful in determining whether a certain set is a manifold or not. Let $M$ be a manifold of dimension $n$, and $S \subseteq M$ be a subset with the subspace topology. Suppose that for all $p \in S$, there exists a chart $(U,\varphi) = (U,x^{1},\ldots,x^{n})$ of $M$ with $p \in U$ and $\varphi(p) = 0 \in \varphi(U) \subseteq \R^{n}$ such that
\begin{align}
    U \cap S = \{(x^{1},x^{2},\ldots,x^{n}) \in U \mid x^{i_{1}} = x^{i_{2}} = \ldots = x^{i_{n-k}} = 0 \text{ for some } 1 \leq i_{1} < i_{2} < \ldots < i_{n-k} \leq n\}
\end{align}
for some fixed integer $0 \leq k \leq n$ (via switching of coordinates, we may assume that the last $n-k$ coordinates are zero). Such a chart $(U,\varphi)$ of $M$ is said to be adapted to $S$. Then, we say that $S$ is a \eax{regular submanifold} of $M$ of dimension $k$ if for all $p \in S$, there exists a chart $(U,\varphi)$ of $M$ with $p \in U$ such that $(U,\varphi)$ is adapted to $S$. Therefore, if $(U,\varphi)$ is adapted to $S$, then $\varphi(U \cap S)$ consists of only those points that look like $(r^{1},\ldots,r^{k},0,\ldots,0) \in \R^{n}$, and hence $\varphi(U \cap S)$ is an open subset of $\R^{k}$. If we define $\varphi_{S}:U \cap S \to \R^{k}$ as $\varphi_{S}(p) = (x^{1}(p),\ldots,x^{k}(p))$, then $\varphi_{S}(U \cap S) = \varphi(U \cap S) \cap \R^{k}$.

\begin{theorem}
    Consider the maximal atlas $\cU = \{(U,\varphi)\}$ of $M$, and let $\cA = \{(U \cap S, \varphi_{S}) \mid (U,\varphi) \in \cU\}$, where $(U,\varphi)$ is adapted to $S$ and $\bigcup \{U \cap S \mid (U,\varphi) \in \cU\} = S$. Then $\cA$ is a smooth atlas on $S$, and hence $S$ is a smooth manifold of dimension $k$.
\end{theorem}
The \eax{codimension} of a regular submanifold $S$ of $M$ is defined as $\codim_{M}S = \dim(M) - \dim(S) = n-k$.

Let $f:M \to \R$ be a smooth function on a smooth manifold $M$ of dimension $n$. Also let $0 \in \R$ be a regular value of $f$ with $f^{-1}(0) \neq \emptyset$. Then, $f^{-1}(0)$ is a regular submanifold of $M$ of dimension $n-1$, or codimension $1$. Note that $0$ here is a red herring; we can replace $0$ with any regular value $c \in \R$ of $f$ with $f^{-1}(c) \neq \emptyset$, and the same result holds. 

For a proof of this, let $p \in f^{-1}(0) \neq \emptyset$ so that $f(p) = 0$. Since $0$ is a regular value of $f$, we have that $f_{\ast,p}:T_{p}M \to T_{0}\R$ is surjective. $T_{p}M$ has a basis $\left\{\frac{\partial}{\partial x^{1}}|_{p},\ldots,\frac{\partial}{\partial x^{n}}|_{p}\right\}$, and $T_{0}\R$ has a basis $\left\{\frac{d}{dt}|_{0}\right\}$. The $1 \times n$ Jacobian matrix of $f$ at $p$ is given by
\begin{align}
    (d f)_{p} = \left( \frac{\partial f}{\partial x^{1}}(p),\ldots,\frac{\partial f}{\partial x^{n}}(p) \right).
\end{align}
Since $f_{\ast,p}:T_{p}M \to T_{0}\R$ is surjective, we have that $\rank(f_{\ast,p}) = 1$, and hence $(d f)_{p} \neq 0$; at least coordinate function $x^{i}$ has $\frac{\partial f}{\partial x^{i}}(p) \neq 0$. Without the loss of generality, we may assume that $\frac{\partial f}{\partial x^{1}}(p) \neq 0$. Then, by the inverse function theorem, there exists an open neighbourhood $V$ of $p$ such that $f:V \to f(V) \subseteq \R$ is a diffeomorphism. Let $(U,\varphi) = (U,x^{1},\ldots,x^{n})$ be a chart on $M$ with $p \in U \subseteq V$. Then, we have
\begin{align}
    U \cap f^{-1}(0) = \{(x^{1},\ldots,x^{n}) \in U \mid x^{1} = 0\},
\end{align}
which shows that $(U,\varphi)$ is adapted to $f^{-1}(0)$. Since $p \in f^{-1}(0)$ was arbitrary, we have that $f^{-1}(0)$ is a regular submanifold of $M$ of dimension $n-1$, or codimension $1$. In general, this result generalizes to the following theorem.

\begin{theorem}
    Let $F:N \to M$ be a smooth map between smooth manifolds of dimension $n$ and $m$, respectivel, with $n \geq m$. Let $q \in M$ be a regular value of $F$ with $F^{-1}(q) \neq \emptyset$. Then, $F^{-1}(q)$ is a regular submanifold of $N$ of dimension $n-m$, or codimension $m$.
\end{theorem}

\begin{proof}
    Let $p \in F^{-1}(q)$; then $F_{\ast,p}:T_{p}N \to T_{q}M$ is surjective. Let $(U,\varphi) = (U,x^{1},\ldots,x^{n})$ be a chart in $N$ such that $\varphi(p) = 0 \in \R^{n}$, and let $(V,\psi) = (V,y^{1},\ldots,y^{m})$ be a chart in $M$ such that $\psi(q) = 0 \in \R^{m}$. We may assume $F(U) \subseteq V$. The matrix of $F_{\ast,p}$ with respect to the bases $\left\{\frac{\partial}{\partial x^{1}}|_{p},\ldots,\frac{\partial}{\partial x^{n}}|_{p}\right\}$ and $\left\{\frac{\partial}{\partial y^{1}}|_{q},\ldots,\frac{\partial}{\partial y^{m}}|_{q}\right\}$ is given by $\left( \frac{\partial F^{i}}{\partial x^{j}}(p) \right)_{1 \leq i \leq m, 1 \leq j \leq n}$; this has rank $m$ since $F_{\ast,p}$ is surjective. Without the loss of generality, we may assume that the $m \times m$ submatrix $\left( \frac{\partial F^{i}}{\partial x^{j}}(p) \right)_{1 \leq i,j \leq m}$ is invertible. Then, by the inverse function theorem, there exists an open neighbourhood $W$ of $p$ such that $F:W \to F(W)$ is a diffeomorphism. Let $(U,\varphi) = (U,x^{1},\ldots,x^{n})$ be a chart on $N$ with $p \in U \subseteq W$. Then, we have
    \begin{align}
        U \cap F^{-1}(q) = \{(x^{1},\ldots,x^{n}) \in U \mid x^{1} = x^{2} = \ldots = x^{m} = 0\},
    \end{align}
    which shows that $(U,\varphi)$ is adapted to $F^{-1}(q)$. Since $p \in F^{-1}(q)$ was arbitrary, we have that $F^{-1}(q)$ is a regular submanifold of $N$ of dimension $n-m$, or codimension $m$.
\end{proof}

\textit{August 19th.}

\begin{example}
    Consider the map $f:\R^{3} \to \R$ defined as $f(x,y,z) = x^{2}+y^{2}+z^{2}$. Then, $S^{2} = f^{-1}(1) = \{(x,y,z) \in \R^{3} \mid x^{2}+y^{2}+z^{2} = 1\}$. We have $f_{\ast,p}:T_{p}\R^{3} \to T_{1}\R = \R$ given by $f_{\ast,p}\left(\frac{\partial}{\partial x^{i}}|_{p}\right) = \sum_{j=1}^{3} \frac{\partial f}{\partial x^{j}}(p) \frac{\partial}{\partial y}|_{1}$, where $y$ is the coordinate function on $\R$. Then, the Jacobian matrix of $f$ at $p$ is given by $(d f)_{p} = \left( \frac{\partial f}{\partial x}(p), \frac{\partial f}{\partial y}(p), \frac{\partial f}{\partial z}(p) \right) = (2x, 2y, 2z)$. The origin is the only critical point of $f$ and is excluded from $S^{2}$. Thus, $1 \in \R$ is a regular value of $f$, and hence $S^{2} = f^{-1}(1)$ is a regular submanifold of $\R^{3}$ of dimension $3-1 = 2$.
\end{example}

\begin{example}
    Consider the determinant map $\det:\R^{n^{2}} \to \R$, interpreting a point $p \in \R^{n^{2}}$ as an $n \times n$ matrix $A \in M_{n}(\R)$. Look at $SL_{n}(\R) = \det^{-1}(1) = \{A \in M_{n}(\R) \mid \det(A) = 1\}$. We have $\det_{\ast,A}:T_{A}\R^{n^{2}} \to T_{1}\R = \R$, where the bases are $\{\frac{\partial}{\partial x_{ij}}\}_{1 \leq i,j \leq n}$ and $\{\frac{d}{dt}|_{1}\}$, respectively. Note that $\frac{\partial \det}{\partial x_{ij}}(A)$ is the $(i,j)$-th cofactor of $A$; not all of these cofactors can be zero since otherwise we would have $\det(A) = 0$ but $A \in \det^{-1}(1)$. Thus, $1 \in \R$ is a regular value of $\det$, and hence $SL_{n}(\R) = \det^{-1}(1)$ is a regular submanifold of $\R^{n^{2}}$ of dimension $n^{2}-1$.
\end{example}


The following is a corollary of the inverse function theorem; it may be generalized to the case of smooth maps between smooth manifolds.

\begin{corollary}
    Let $p \in \R^{n}$ have an open neighbourhood $U \subseteq \R^{n}$, and let $f:U \to \R^{m}$ be a smooth map with $f = (f_{1},\ldots,f_{m})$. Suppose $f$ has rank $k$ in a neighbourhood of $p$ in $U$; that is, $\left(\frac{\partial f_{i}}{\partial x_{j}}(q)\right)$ has rank $k$ for all $q$ in this possibly smaller neighbourhood of $p$. Then there exists a diffeomorphism $G:V \to G(V)$, with $G(V)$ open in $\R^{n}$ and $p \in V$ open in $U$, and a diffeomorphism $F:W \to F(W)$, with $F(W)$ open in $\R^{m}$ and $f(p) \in W$ open in $\R^{m}$, such that $f(V) \subseteq W$, and $(FfG^{-1})(x^{1},\ldots,x^{n}) = (x^{1},\ldots,x^{k},0,\ldots,0)$ for all $(x^{1},\ldots,x^{n}) \in G(V)$.    
\end{corollary}

\begin{proof}
    Let $p \in \R^{n}$ have an open neighbourhood $U \subseteq \R^{n}$, and let $f:U \to \R^{m}$ be a smooth map with $f = (f_{1},\ldots,f_{m})$. Suppose $f$ has rank $k$ in a neighbourhood of $p$ in $U$. Then, without the loss of generality, we may assume that the $k \times k$ submatrix $\left(\frac{\partial f_{i}}{\partial x_{j}}(q)\right)_{1 \leq i,j \leq k}$ is invertible for all $q$ in this possibly smaller neighbourhood of $p$. Define the smooth map $F:U \to \R^{n}$ as $F(x^{1},\ldots,x^{n}) = (f_{1}(x^{1},\ldots,x^{n}),\ldots,f_{k}(x^{1},\ldots,x^{n}),x^{k+1},\ldots,x^{n})$. Then, the Jacobian matrix of $F$ at $q \in U$ is given by
    \begin{align}
        (dF)_{q} = \begin{pmatrix}
            \frac{\partial f_{1}}{\partial x_{1}}(q) & \cdots & \frac{\partial f_{1}}{\partial x_{k}}(q) & \frac{\partial f_{1}}{\partial x_{k+1}}(q) & \cdots & \frac{\partial f_{1}}{\partial x_{n}}(q) \\
            \vdots & \ddots & \vdots & \vdots & \ddots & \vdots \\
            \frac{\partial f_{k}}{\partial x_{1}}(q) & \cdots & \frac{\partial f_{k}}{\partial x_{k}}(q) & \frac{\partial f_{k}}{\partial x_{k+1}}(q) & \cdots & \frac{\partial f_{k}}{\partial x_{n}}(q) \\
            0 & 0 & 0 & 1 & 0 & 0\\
            0 & 0 & 0 & 0 & 1 & 0\\
            0 & 0 & 0 & 0 & 0 & 1
        \end{pmatrix}.
    \end{align}
    Thus, by the inverse function theorem, there exists an open neighbourhood $V$ of $p$ with $V \subseteq U$ such that $F:V \to F(V)$ is a diffeomorphism. Consider the map $F^{-1}:F(V) \to V$; a point $(x^{1},\ldots,x^{n}) \in F(V)$ is mapped to $(x^{1},\ldots,x^{k},f^{k+1},\ldots,f^{m})$. The Jacobian matrix of $F^{-1}$ at $(x^{1},\ldots,x^{n}) \in F(V)$ is given by
    \begin{align}
        (dF^{-1})_{(x^{1},\ldots,x^{n})} = \begin{pmatrix}
            1 & 0 & 0 & 0 & 0 & 0 \\
            0 & 1 & 0 & 0 & 0 & 0 \\
            0 & 0 & 1 & 0 & 0 & 0 \\
            \ast & \ast & \ast & \frac{\partial f^{k+1}}{\partial x^{k+1}} & \cdots & \frac{\partial f^{k+1}}{\partial x^{n}} \\
            \ast & \ast & \ast & \vdots & \ddots & \vdots \\
            \ast & \ast & \ast & \frac{\partial f^{m}}{\partial x^{k+1}} & \cdots & \frac{\partial f^{m}}{\partial x^{n}}
        \end{pmatrix}.
    \end{align}
    We know this matrix to have rank $k$ by hypothesis. Thus, the $(m-k) \times (n-k)$ submatrix $\left(\frac{\partial f^{i}}{\partial x^{j}}(q)\right)_{k+1 \leq i \leq m, k+1 \leq j \leq n}$ must have rank $0$; that is, $\frac{\partial f^{i}}{\partial x^{j}}(q) = 0$ for all $k+1 \leq i \leq m$ and $k+1 \leq j \leq n$. The functions $f^{k+1},\ldots,f^{m}$ are then functions of only $x^{1},\ldots,x^{k}$. Define the smooth map $G:f(V) \to \R^{m}$ as $(y^{1},\ldots,y^{m}) = (y^{1},\ldots,y^{k},y^{k+1}-f^{k+1},\ldots,y^{m}-f^{m})$. Then, the Jacobian matrix of $G$ at $(y^{1},\ldots,y^{m}) \in f(V)$ is given by
    \begin{align}
        (dG)_{(y^{1},\ldots,y^{m})} = \begin{pmatrix}
            1 & 0 & 0 & 0 & 0 & 0 \\
            0 & 1 & 0 & 0 & 0 & 0 \\
            0 & 0 & 1 & 0 & 0 & 0 \\
            -\frac{\partial f^{k+1}}{\partial y^{1}} & \cdots & -\frac{\partial f^{k+1}}{\partial y^{k}} & 1 & 0 & 0 \\
            \vdots & \ddots & \vdots & 0 & 1 & 0 \\
            -\frac{\partial f^{m}}{\partial y^{1}} & \cdots & -\frac{\partial f^{m}}{\partial y^{k}} & 0 & 0 & 1
        \end{pmatrix}.
    \end{align}
    This matrix is of full rank, so $G$ is a diffeomorphism from $f(V)$ to $G(f(V))$. Then, we have $(GF)(x^{1},\ldots,x^{n}) = (f_{1}(x^{1},\ldots,x^{n}),\ldots,f_{k}(x^{1},\ldots,x^{n}),0,\ldots,0)$ for all $(x^{1},\ldots,x^{n}) \in V$. Thus, we have found diffeomorphisms $G$ and $F$ such that $(GF)(x^{1},\ldots,x^{n}) = (x^{1},\ldots,x^{k},0,\ldots,0)$ for all $(x^{1},\ldots,x^{n}) \in G(V)$.
\end{proof}

The generalization is as follows. Let $f:N \to M$ be a smooth map between smooth manifolds of dimension $n$ and $m$, respectively. Let $p \in N$ be a point such that $f$ has a rank $k$ in a neighbourhood of $p$ in $N$. Then there exists a chart $(U,\varphi)$ centred at $p$ in $N$ and a chart $(V,\psi)$ centred at $f(p)$ in $M$ such that $f(U) \subseteq V$, and $\psi \circ f \circ \varphi^{-1}(x^{1},\ldots,x^{n}) = (x^{1},\ldots,x^{k},0,\ldots,0)$ for all $(x^{1},\ldots,x^{n}) \in \varphi(U)$.
\\

\textit{August 24th.}
\begin{theorem}
    Let $f:N \to M$ be a smooth map between smooth manifolds of dimension $n$ and $m$, respectively. Let $c \in M$ be such that $f$ has a rank $k$ in a neighbourhood of $f^{-1}(c) \neq \emptyset$ in $N$. Then $f^{-1}(c)$ is a regular submanifold of $N$ of dimension $n-k$.
\end{theorem}

% \begin{proof}
%     We wish to find smooth compatible charts on open subsets of $S = f^{-1}(c)$ which will prove that $S$ is a regular submanifold of $N$. Let $p \in S$ so that $f(p) = c$. Then, by the next theorem, there exists a chart $(U,\varphi)$ centred at $p$ in $N$ and a chart $(V,\psi)$ centred at $c$ in $M$ such that $\psi f \varphi^{-1}(x^{1},\ldots,x^{n}) = (x^{1},\ldots,x^{k},0,\ldots,0) \in \R^{m}$ for all $(x^{1},\ldots,x^{n}) \in \varphi(U)$. Then, with $\psi(c) = \psi(f(p)) = 0 \in \R^{m}$, we have $S \cap U = (\psi f)^{-1}(0) \cap U$, and $\varphi(S \cap U) = \varphi((\psi f)^{-1}(0) \cap U) = \varphi((\psi f)^{-1}(0))$
% \end{proof}

\begin{example}
    Look at the map $f:GL_{n}(\R) \to GL_{n}(\R)$ define by $f(A) = A^{t}A$. Then, $O_{n}(\R) = f^{-1}(I_{n}) = \{A \in GL_{n}(\R) \mid A^{t}A = I_{n}\}$. For now, the dimension of $O_{n}(\R)$ will remain unknown, but we can still show that it is a regular submanifold of $GL_{n}(\R)$. Let $C \in GL_{n}(\R)$. Then $f(AC) = C^{t}f(A)C$. Note that conjugation by $C$ is a diffeomorphism from $GL_{n}(\R)$ to itself, and hence $f$ has the same rank at $A$ and $AC$. We shall show that the ranks of $f_{\ast,A}$ and $f_{\ast,B}$ are the same for all $A,B \in GL_{n}(\R)$. We can break down $B = AC$ for some $C \in GL_{n}(\R)$. Then $f \circ r_{C} = l_{C^{t}} \circ r_{c} \circ f$, where $r_{C}:GL_{n}(\R) \to GL_{n}(\R)$ is right multiplication by $C$ and $l_{C^{t}}:GL_{n}(\R) \to GL_{n}(\R)$ is left multiplication by $C^{t}$. Then, we have
    \begin{align}
        f_{\ast,AC} (r_{C})_{\ast,A} = (l_{C^{t}})_{\ast,A^{t}AC}(r_{C})_{\ast,A^{t}A}f_{\ast,A}.
    \end{align}
    The rank of $f_{\ast,A}$ is the same as the rank of $f_{\ast,AC}$ since $(r_{C})_{\ast,A}$ and $(l_{C^{t}})_{\ast,A^{t}AC}(r_{C})_{\ast,A^{t}A}$ are both isomorphisms. Thus, $f$ has the same rank at all points in $GL_{n}(\R)$. 
\end{example}


The previous theorem may be generalized to the following theorem.

\begin{theorem}
    Let $f:N \to M$ be a smooth map between smooth manifolds of dimension $n$ and $m$, respectively. Also let $S \subseteq M$ be a regular submanifold of codimension $k$, or dimension $m-k$. Suppose $f_{\ast,p}(T_{p}N) + T_{f(p)}S = T_{f(p)}M$ for all $p \in f^{-1}(S)$. Then, $f^{-1}(S)$ is a regular submanifold of $N$ of codimension $k$, or dimension $n-k$.
\end{theorem}


\begin{theorem}
    Let $f:N \to M$ be a smooth map between smooth manifolds of dimension $n$ and $m$, respectively. Let $p \in N$.
    \begin{enumerate}
        \item Suppose $\rank(f_{\ast,p}) = m$. Then there exists a local chart $(U,\varphi)$ centred at $p$ in $N$ and a local chart $(V,\psi)$ centred at $f(p)$ in $M$ such that $\psi f \varphi^{-1}(x^{1},\ldots,x^{n}) = (x^{1},\ldots,x^{m})$ for all $(x^{1},\ldots,x^{n}) \in \varphi(U)$.
        \item Suppose $\rank(f_{\ast,p}) = n$. Then there exists a local chart $(U,\varphi)$ centred at $p$ in $N$ and a local chart $(V,\psi)$ centred at $f(p)$ in $M$ such that $\psi f \varphi^{-1}(x^{1},\ldots,x^{n}) = (x^{1},\ldots,x^{n},0,\ldots,0)$ for all $(x^{1},\ldots,x^{n}) \in \varphi(U)$.
    \end{enumerate}
\end{theorem}

